\chapter*{勘误表}
\addcontentsline{toc}{chapter}{勘误表}
\markboth{勘误表}{勘误表}

本表汇总中译全过程发现的原书排印疑误，共 325 条，按章分组。正文相应位置均以
“译者注”脚注说明，本表撮其要，供对照原书阅读时参考。条目位置给出译稿小节号与
原书印刷页码（同一小节跨页时给出页码范围）。处置凡分两类：\textbf{径改}指译文
已按酌定改排补正；\textbf{存疑照排}指证据不足或两读可通，译文一仍原书之旧，
仅备案待考。

\subsection*{第\,2\,章}
\begin{description}[font=\normalfont,leftmargin=6em,itemsep=0.2em]
\item[§2.4，26-29] 原书在局部开集 $\Omega$ 的表述中 直接称 $u^{-1}(p)$ 为有限集；定理 2.3.5 在这里给出的是局部孤立性，而有限性还须使用 全局定义域 $\Sigma$ 的紧致性。（径改）
\item[§2.4，26-29] 原书此处仅由 $z_\nu\ne z_0$ 直接推出 $u_0(z_\nu)\ne u_0(z_0)$；这一推论对一般映射并不成立。这里利用 $du_0(z_0)\ne0$ 以及引理 2.4.2 给出的局部单射性补足此步。（径改）
\item[§2.6，35-37] 原书把 $v$ 的靶空间印作 $Q$；然而这里须计算 $v$ 与 $Q$ 的交点，且 $v$ 是保持边界值的 $u:B_1\to M$ 的扰动，故改为 $M$。（径改）
\item[§2.6，35-37] 原书此处把两条分量重合时的 $2\delta$ 写成 有限值；但按本节在定义~\ref{def:2-6-2} 后给出的 $\delta(u)$ 定义及紧接的说明， 非单纯曲线的 $\delta(u)$ 为无穷，与此不相容。译文照录原式，不自行重构作者意图。（存疑照排）
\end{description}

\subsection*{第\,3\,章}
\begin{description}[font=\normalfont,leftmargin=6em,itemsep=0.2em]
\item[§3.1，39-48] 原书此处写作 “\textit{simple $J$-holomorphic spheres}”；但本节随即固定任意紧黎曼曲面 $\Sigma$，且定理~\ref{thm:3-1-6} 与定理~\ref{thm:3-1-8} 也在这一范围内陈述， 故改作“曲线”。（径改）
\item[§3.1，39-48] 原书称 $\mathcal B^{k,p}$ 是 $W^{k,p}(\Sigma,M)$ 的一个连通分支； 但固定同调类一般不能唯一确定映射空间的同伦分支，故改作“若干连通分支之并”。（径改）
\item[§3.3，56-61] 原书此处写作“$(TM,J)$ 的每个线子丛 $L$”；然而该丛的底空间是 $M$，而这里的算子及其分裂定义在 $\mathbb{CP}^1$ 上。由紧邻的分裂式和下文 Dolbeault 算子的定义域可知，应为 $u^*TM$ 的全纯直和项。（径改）
\item[§3.3，56-61] 原书此处及证明末句均断言“当且仅当”。 按 $D$-不变子丛的通常含义，充分性成立，但必要性一般不成立：例如 $\mathbb{CP}^1$ 上有非平凡正合列 $0\to\mathcal O(-2)\to\mathcal O(-1)\oplus\mathcal O(-1) \to\mathcal O\to0$；中项的 Dolbeault 算子满射，而滤过的第一个商的陈数为 $-2$。 这里保留原书断言；下文引理~\ref{lem:3-3-3} 的滤过中，第一个子丛上的限制算子本身 满射，此时由相应的核余核正合列可知，总算子满射当且仅当商算子满射，故该应用不受此 反例影响。（存疑照排）
\item[§3.3，56-61] 原书此处写作 $\widetilde u^*TM$；但 $\widetilde u$ 的靶是 $\widetilde M=S^2\times M$，且紧随其后的 分解包含来自 $S^2$ 因子的陈数为 $2$ 的切向线丛，故改为 $\widetilde u^*T\widetilde M$。（径改）
\item[§3.4，61-68] 原书此处把测试截面写作 $\eta\in W^{1,p}(\Sigma,u^*TM)$；但 $\eta$ 已被假设为取值于 $\Lambda^{0,1}\otimes_Ju^*TM$ 的 $L^q$ 形式，而紧随其后的式子使用的是 $D_u\xi$。故依式~\eqref{eq:3-4-2} 唯一改为 $\xi$。（径改）
\item[§3.4，61-68] 原书此式右端漏印因子 $1/2$。由本证明前一显示式以及 $\eta=\Phi\eta_0$ 可唯一确定该因子；后文只使用这个表达式等于零，故原书后续论证 不受影响。（径改）
\item[§3.5，68-73] 原书此处把映射误写作 $\mathcal F_u(\xi):\mathcal X_u\to\mathcal Y_u$；但 $\mathcal F_u(\xi)$ 是 $\mathcal Y_u$ 中的取值，而紧邻上文定义的映射是 $\mathcal F_u$，故据源靶类型改正。（径改）
\item[§3.5，68-73] 原书在 $\Psi_x$ 的点态定义及紧随其后的说明中把 $\Phi_x$ 两次印作 $\Phi_u$；但此处给定的是 $x\in M$，相关测地线以 $x$ 为基点，且源靶均以 $x$ 标记， 故据类型以及随后从逐点 $\Psi_x$ 得到 $\Psi_u$ 的记号唯一改为 $\Phi_x$。（径改）
\item[§3.5，68-73] 原书把这个映射的定义域印作 $S^2$；但本节从始至终固定的是任意闭黎曼曲面 $\Sigma$，且 $\exp_u(\xi)$ 是逐点指数映射所得的 $\Sigma\to M$ 映射，故改为 $\Sigma$。（径改）
\item[§3.5，68-73] 原书在这个索伯列夫嵌入估计中漏印了 $\xi'$ 的撇号，因而把前一项 $\xi$ 的估计重复了一次。上一行含有因子 $\lVert\xi'\rVert_{L^\infty}$，而紧随其后的 式~\eqref{eq:3-5-4} 要把它估计为 $\lVert\xi'\rVert_{W^{1,p}}$；常数 $c_3=(c_0+1)c_0^2c_1c_2$ 也印证了这个唯一补正。（径改）
\end{description}

\subsection*{第\,4\,章}
\begin{description}[font=\normalfont,leftmargin=6em,itemsep=0.2em]
\item[§4.1，75-80] 原书作“Schwartz Lemma”；这里所指是共形映射的 Schwarz 引理，故据标准名称改正。（径改）
\item[§4.1，75-80] 原书正文拼作“Bourgeous”；同书参考文献 \srcite{46} 的作者 条目作“F. Bourgeois”，据此改正。（径改）
\item[§4.1，75-80] 原书此处写作 “$u:\mathbb D\to M$”；但这里正在证明定理~\ref{thm:4-1-2}\textup{(i)} 中定义在 穿孔圆盘上的映射可连续延拓，且下文考察 $z\to0$ 的极限，故补回 $\setminus\{0\}$。（径改）
\item[§4.1，75-80] 原书此处写作 $E(u)<\delta$；但按本书约定 $E(u)=\frac12\int|du|^2$，而上述估计的假设是 $\int|du|^2<\delta$，故充分条件应为 $E(u)<\delta/2$。等价地也可重命名阈值；命题只需某个正下界，结论不受影响。（径改）
\item[§4.1，75-80] 原书下式第二个下确界的映射域写作 $(B,\partial B)$，但 $B$ 在此尚未定义， 且命题~\ref{prop:4-1-4} 明确以闭单位圆盘 $\mathbb D$ 为域；故将这个定义域符号补正为 $(\mathbb D,\partial\mathbb D)$。（径改）
\item[§4.2，80-85] 原书此处写作 $z^\nu\in\Omega$；但此前 $z^\nu\in\Sigma$，而坐标域 $\Omega\subset\mathbb C$，故据坐标卡 $\phi:\Omega\to\Sigma$ 的源靶改作 $z^\nu\in\phi(\Omega)$。（径改）
\item[§4.2，80-85] 原书最后一行写作 $\le\lim_{\nu\to\infty}$，但并未建立右侧 能量列的极限存在；这里的目标式~\eqref{eq:4-2-1} 使用 $\liminf$，且由逐项包含关系 唯一得到的也是 $\liminf$，故据此补正。（径改）
\item[§4.3，85-91] 原书此处写作 “and $L$ is totally real”。但 $L$ 为全实子流形已是引理~\ref{lem:4-3-1} 的假设， 与度量选取无关；紧接着所引的引理~\ref{lem:4-3-4} 构造的正是使 $L$ 全测地的度量， 故据此改正。（径改）
\item[§4.3，85-91] 原书在 $a\ge0$ 的同时直接写出 $\int_{B_r}w<\pi/(8a)$，但该商式在 $a=0$ 时无定义；证明的步骤 3 明确把 $a=0$ 作为无须积分小量条件的次调和情形处理， 故按书内论证补明这一约定。（径改）
\item[§4.3，85-91] 原书此处写作 “$0\le\rho\le\varepsilon$”；但紧随其后的不等式含有 $1/\rho^2$，在 $\rho=0$ 处没有定义，故将左端点改为严格不等号。（径改）
\item[§4.3，85-91] 原书此处写作“由有限多个 截断函数组成的单位分解”。引理并未假设 $M$ 或 $L$ 紧；一般情形所需且足够的是局部有限 单位分解，故据此改正。（径改）
\item[§4.3，85-91] 原书此处写作 “$\Delta w$ 满足引理 4.3.3 的假设”； 但引理~\ref{lem:4-3-3} 的非负性、积分小量与微分不等式均以函数 $w$ 为对象，故改作 “$w$ 满足”。（径改）
\item[§4.4，91-96] 原书此处写作 $|a(u_\gamma)|$；但 $a$ 在本节中只定义在边界闭路 $\gamma$ 上，而 $u_\gamma$ 是填充 圆盘，故据作用量的参数类型改作 $|a(\gamma)|$。（径改）
\item[§4.4，91-96] 原书此处写作 $\alpha(\theta+2\pi)=\alpha(\theta)$；这会使 $\alpha$ 不可能成为圆周的提升微分同胚。由闭路参数的周期性，唯一应改为 $\alpha(\theta+2\pi)=\alpha(\theta)+2\pi$。（径改）
\item[§4.4，91-96] 原书断言它对任意 $Q\in\operatorname{GL}(n,\mathbb R)$ 都被 $\omega$ 驯服，但一般并不成立，例如 $Q=-I$ 即给出反例；相应二次型为正定才有驯服性， 等价地须要求 $Q$ 的对称部分正定。下述具体矩阵族在 $0<\varepsilon\le1$ 时对称部分为 $\varepsilon I$，故后续例族不受这一一般性过强的断言影响。（径改）
\item[§4.6，99-105] 原书此处的能量链末项写作 $\liminf_\nu E(u^\nu;B_\varepsilon(z_0))$；但 $z_0$ 允许位于边界，原映射只定义在 $\mathbb H$ 中，而上面的包含关系与目标式~\eqref{eq:4-6-1} 均含 $B_\varepsilon(z_0)\cap\mathbb H$，故补上 $\cap\mathbb H$。（径改）
\item[§4.6，99-105] 原书此处写作 $E(v;B_R)>\alpha\hbar$；但本情形中的 $v$ 只定义在 $\mathbb H$ 上，且紧接着的能量链 也使用 $B_R\cap\mathbb H$，故据定义域唯一补上 $\cap\mathbb H$。（径改）
\item[§4.7，105-114] 原书此处写作“满足该引理的假设”（the assumptions of the lemma）； 但引理~\ref{lem:4-7-5} 所需的重标度形式和无能量损失条件分别到步骤 2、3 才建立， 此处坐标变换保持的是命题~\ref{prop:4-7-2} 的假设，故改正引用对象。（径改）
\item[§4.7，105-114] 原书此处固定 $\mu=1/2$；但定理~\ref{thm:4-4-1}(iii) 在一般驯服情形 只给某个未受控制的等周常数，而引理~\ref{lem:4-7-3} 证明中的 $\mu$ 必须据此选取， 故不能一般地指定为 $1/2$。这里沿用该引理及注~\ref{rem:4-7-4} 给出的某个固定 $0<\mu<1$；原书步骤 3 的 $\delta/R$ 是 $\mu=1/2$ 时的特例，一般情形的对应项仍趋于零， 所以反证不变。（径改）
\end{description}

\subsection*{第\,5\,章}
\begin{description}[font=\normalfont,leftmargin=6em,itemsep=0.2em]
\item[§5.1，115-122] 原书把“建立格罗莫夫极限的唯一性”与“把收敛推广到稳定映射列”两项都归于 第 5.4 节；但第 5.4 节只证明唯一性，而稳定映射列的格罗莫夫收敛到第 5.5 节开头才定义， 故分别改作第 5.4 节与第 5.5 节。（径改）
\item[§5.1，115-122] 原书此处写成压缩只有一或两个“标记点”的顶点；但定义~\ref{def:5-1-1} 中的稳定性由特殊点数决定，而特殊点包括节点与标记点，故改作“特殊点”。（径改）
\item[§5.2，122-126] 原书此处写作 $u(z_0)=v(\infty)$；但本段已定义 $z_\infty:=z_{\beta\alpha}$，命题~\ref{prop:4-7-2} 的坐标不变结论及紧接着的 $u_\beta(z_{\beta\alpha})$ 都唯一要求此处为 $v(z_\infty)$。（径改）
\item[§5.2，122-126] 原书此处写作 $\psi^\nu:=\phi_{\beta\alpha}^\nu$；但按 $\phi_{\alpha\beta}^\nu=(\phi_\alpha^\nu)^{-1}\circ\phi_\beta^\nu$，只有取 $\psi^\nu=\phi_{\alpha\beta}^\nu$ 才同时满足 $u_\beta^\nu=u_\alpha^\nu\circ\psi^\nu$，并使下文 $(\psi^\nu)^{-1}((\phi_\alpha^\nu)^{-1}(z)) =(\phi_\beta^\nu)^{-1}(z)$。原书的反向次序不能推出所需距离估计。（径改）
\item[§5.2，122-126] 原书此处把节点写成不存在的 $u_{\alpha\beta}$，从而出现 $u_\alpha(u_{\alpha\beta})$；由本段定义 $z_0:=z_{\alpha\beta}$ 以及 命题~\ref{prop:4-7-2} 的结论，唯一有类型的表达式是 $u_\alpha(z_{\alpha\beta})$。（径改）
\item[§5.3，126-134] 原书只写“只要 $j_k=j_{k'}$，就有 $z_k\ne z_{k'}$”，遗漏了 $k\ne k'$；按字面取 $k=k'$ 会自相矛盾，而这里所需且由后文唯一确定的条件，是同一 母顶点上的不同子分支具有互异的附着点，故补上 $k\ne k'$。（径改）
\item[§5.3，126-134] 原书此处写作 $i\ge1$；但根分支 $i=1$ 没有指向母分支的节点 $\infty$，而且前文已经规定 $u_1$ 非常值，故这里唯一相关的指标范围是 $i\ge2$。（径改）
\item[§5.3，126-134] 原书称 $\hbar$ 为这些非零权重的“最小值”；命题~\ref{prop:4-1-4} 只保证每个非零权重至少为 $\hbar$，下文也只需要这一统一正下界，故据此改正。（径改）
\item[§5.3，126-134] 原书在此处及随后引理~\ref{lem:5-3-4} 的假设中，两次把集合下标的 起点印作未定义的 $i$；此时已有且新增点必须避开的是全部 $k$ 个标记点，故统一改为 $1$。（径改）
\item[§5.3，126-134] 原书此处印作 $\zeta_{i+1}$；但式~\eqref{eq:5-3-5} 只定义了按树顶点 $\alpha$ 标记的量 $\zeta_\alpha$，故补为 $\zeta_{\alpha_{i+1}}$。（径改）
\item[§5.4，134-139] 原书此处以及步骤 4 都把标记点指标的上界印作 $k$；但 定理~\ref{thm:5-4-1} 从头到尾固定的是 $n$ 个标记点，而本证明没有定义 $k$，故两处 统一改为 $n$。（径改）
\item[§5.4，134-139] 原书在此未写 $\alpha\ne\beta$；但上文只对不同顶点定义了 $z_{\alpha\beta}$，而边关系也只连接不同顶点，故补上这一条件。（径改）
\item[§5.5，139-147] 原书此处写作 $\rho_\varepsilon(\mathbf x,\mathbf x')>0$；但本节定义允许 其取零，且注~\ref{rem:5-5-7} 明确指出等价稳定映射之间的该量为零，故补为非负。（径改）
\item[§5.5，139-147] 原书把极限稳定映射所属空间印作 $\mathcal{SC}_{0,n}(M;J^\nu)$；但 $J^\nu$ 随 $\nu$ 变化，而格罗莫夫极限按 定义~\ref{def:5-5-1} 是 $J$-全纯稳定映射，故补为 $\mathcal{SC}_{0,n}(M;J)$。（径改）
\item[§5.5，139-147] 原书称命题~\ref{prop:4-1-4} 的最大阈值 $\hbar(M,\omega,J)$ 连续依赖于 $J$，并据此推出 $\inf_\nu\hbar(M,\omega,J^\nu)>0$；正文并未证明这个由可实现全纯球面能量的下确界 所定义函数的连续性，一般也不应期待这种连续性。此处只需局部一致的正下界，而 $J^\nu\to J$ 所成紧参数族上的平均值估计常数可一致选取，已足以得到该结论。（径改）
\item[§5.5，139-147] 原书未单独处理树同态不塌缩首条边时可能出现的 $\ell=0$，而仍写出 $z_{\gamma_0\gamma_\ell}$；这会成为未定义的 $z_{\gamma_0\gamma_0}$。此时（节点）公理 直接给出所需极限；以下原证明则适用于 $\ell\ge1$，引理结论不变。（径改）
\item[§5.5，139-147] 原书把上述点值极限的中间项印作 $(\phi_\alpha^\nu)^{-1}\circ\phi_{\alpha_i}^\nu(B_\delta(z_i))$，这是一个集合， 与等式两侧的点类型不合。令 $\zeta^\nu=(\phi_{\alpha_i}^\nu)^{-1}(z_i^\nu)\to z_i$，再用 (i) 在 $z_i\ne z_{\alpha_i\alpha}$ 的邻域上一致收敛，即唯一得到正文中的点值等式。（径改）
\item[§5.6，147-151] 原书断言拓扑“唯一”，但随后的证明只构造并验证了序列化拓扑 $\mathcal U(C)$，即证明 $C(\mathcal U(C))=C$；一般拓扑未必由其收敛序列唯一决定。若把 范围限定为由序列刻画的序列型拓扑，则可得到相应的唯一性，但原文未加这一限制。（径改）
\item[§5.6，147-151] 原书此处由第一可数性与可数稠密子集直接推出第二可数性；一般而言，这一 推理并不成立。要得到所声称的可数基，还需另行证明该格罗莫夫拓扑第二可数。（径改）
\end{description}

\subsection*{第\,6\,章}
\begin{description}[font=\normalfont,leftmargin=6em,itemsep=0.2em]
\item[§6.1，153-158] 原书把命题的输入空间印作带星号的 $\widetilde{\mathcal M}^{*}_{0,T}$；但本节刚把星号空间定义为简单稳定映射空间，而命题前文、 结论及以下三步证明均从任意稳定映射出发，故去掉星号。（径改）
\item[§6.2，158-164] 原书在此处及第~\ref{sec:6-3} 节开头两次把原像映射误印为 $\pi^E$；本章定义的相应映射是边求值映射 $\operatorname{ev}^E$，且原像条件正是 节点两端取值相等，故两处均改正。（径改）
\item[§6.2，158-164] 原书此处把作用对象印作 $\widetilde{\mathcal M}^{*}(\{A_\alpha\};J)$，漏掉下标 $0,T$；前一显示刚定义的原像空间 才含有节点数据、其商才是本定理中的稳定映射模空间，故补回下标。（径改）
\item[§6.2，158-164] 原书把 $\Delta^E$ 两度描述为这些子流形的“并”，并进而称其为由不同维数 子流形组成的分层空间；但定义~\ref{def:6-2-1} 要求同时满足所有边等式，且下文使用的 $\operatorname{codim}\Delta^E=2ne(T)$ 也只与交（等价于成对子对角线的乘积）一致，故据此 改正。（径改）
\item[§6.2，158-164] 原书把交集指标印作所有实数 $K>0$，随即称其为“可数交”。这里的 $K$-有界子集随 $K$ 单调并穷竭模空间，故取正整数这一可数共尾子集即可，依此改正。（径改）
\item[§6.3，165-168] 原书此处把模空间的同调类参数印作总类 $A$；但定理固定的是整族 $\{A_\alpha\}$，本节前后的普适模空间也都以这族分支类为参数，故改正。（径改）
\item[§6.3，165-168] 原书把这里的向量场误印作带两个顶点下标的 $\xi_{\alpha_i\beta}$；切向量元组只定义了按顶点标号的 $\xi_\alpha$，故改为 $\xi_{\alpha_i}$。（径改）
\item[§6.4，168-171] 原书脚注 1 把“反典范丛”误写成“典范丛”；但正文刚定义 $K^*=\Lambda^{n,0}TM$，且脚注随后要求 凯勒 类等于 $c_1(TM)$，两者都唯一确定此处 应为反典范丛。（径改）
\item[§6.4，168-171] 原书把该球面的陈数写成 $2-m$；但 $[u_m]=(1,-m)$，而 $c_1(T(S^2\times S^2))=(2,2)$，故其陈数唯一为 $2-2m$。 同句又误称“Exercise 3.3.6”，实际对应例~\ref{ex:3-3-6}，一并改正。（径改）
\item[§6.5，171-177] 原书括号内断言只由维数条件自动得到 $\dim\Omega_f\le d-2$，却没有保证定义~\ref{def:6-5-1} 还要求的 $\overline{f(V)}$ 紧；例如恒等映射 $\mathbb R^d\to\mathbb R^d$ 的像闭包并不紧。若另假设 $X$ 紧或 $f(V)$ 预紧， 该断言才成立。原书在例~\ref{ex:6-5-12} 又重复了同一缺条件说法。（径改）
\item[§6.5，171-177] 原书直接称每个这样的二维复形本身都可赋予闭光滑曲面结构；一般复形的 一个顶点可能有多个圆周状局部分支，须先分离这些分支（即作顶点正规化）。这一处理不改变 所表示的二维循环，并给出正文中的闭曲面。（径改）
\item[§6.5，171-177] 原书在双向断言中只假设 $Y$ 是 $Z$ 的子流形；逆向还须假设 $Y$ 在 $Z$ 中闭（等价地，此处的包含映射恰当），才能由像在 $Z$ 中有紧闭包推出其在 $Y$ 中 也有紧闭包。例如，$Y=S^2\setminus\{p\}\subset S^2=Z$ 且 $e=\operatorname{id}_Y$ 时，复合映射是二维伪循环，而 $e:Y\to Y$ 的像闭包并不紧。后文只用正向推论，不受影响。（径改）
\item[§6.5，171-177] 这里再次沿用了例~\ref{ex:6-5-12} 前的缺条件说法；要断言任意 $\pi,e$ 都是伪循环，仍须假设各自的像有紧闭包（例如 $X,Y$ 紧）。后面的具体反例取 $X=Y=S^2$，不受此缺口影响。（径改）
\item[§6.5，171-177] 与习题~\ref{ex:6-5-2}(ii) 的括号说明相同，这一断言还需要 $\overline{f(U)}$ 紧；单有维数不等式并不足够。此处保留原书表述，疑点见前注。（存疑照排）
\item[§6.6，177-181] 原书此处写成“至少两条边”；但一条边的稳定映射同样是主层的边界， 而下文的维数论证明确对所有 $e(T)>0$ 的树进行。故唯一补为“至少一条边”。（径改）
\item[§6.7，182-199] 原书把末项印作 $2(k-2\#I)$；但解空间指标为 $2n+2c_1(A)$，加入 $k$ 个可变点贡献 $2k$ 维，而固定每个定义域点只减 $2$ 维。 定理~\ref{thm:6-7-11}\textup{(i)} 在 $e(T)=0$ 时以及本节稍后的同伦模空间维数也都给出 $2n+2c_1(A)+2k-2\#I$，故据此补正。（径改）
\item[§6.7，182-199] 若这里的“没有自同构”意指没有非平凡自同构，则原书的范围 $g>1$ 并不 包含亏格 $2$：每个亏格 $2$ 曲面都有超椭圆对合；相应的通有无非平凡自同构结论应取 $g>2$。若这里只意指没有连续的或无穷小的自同构，则原范围成立。（径改）
\item[§6.7，182-199] 原书在式~\eqref{eq:6-7-6} 中把体积形式的下标印作 $S^2$，但本引理的定义域 是任意闭黎曼曲面 $\Sigma$，紧随其后的局部计算也使用 $\mathrm{dvol}_\Sigma$。原书随后又把 共形度量 $\lambda^2(ds^2+dt^2)$ 误称为体积形式；正文分别改为正确的 $\mathrm{dvol}_\Sigma$，并明确相应体积形式为 $\lambda^2ds\wedge dt$。（径改）
\item[§6.7，182-199] 原书把椭圆正则性提升后的空间印作 $W^{\ell,p}$；但 $\eta$ 起初属于 $L^q$，其中 $q$ 是 $p$ 的对偶指数，作用于伴随算子 $D_u^*$ 的椭圆正则性保持这一积分 指数，故应为 $W^{\ell,q}$。（径改）
\item[§6.7，182-199] 原书在这里把参数区间印作 $0\le\lambda<1$；但同伦的定义域是闭区间 $[0,1]$，且命题~\ref{prop:6-7-9}\textup{(iii)} 随即把 $\lambda=1$ 的模空间列为配边的一侧， 故补上端点 $\lambda=1$。（径改）
\item[§6.7，182-199] 原书把式~\eqref{eq:6-7-9} 左端的定义域印作 $\Sigma$；但 $v:S^2\to M$，式~\eqref{eq:6-7-8} 中 $D_{z,v}$ 的靶空间以及右端 $dv\circ j_{S^2}$ 都定义在 $S^2$ 上，故补正为 $S^2$。（径改）
\item[§6.7，182-199] 原书的条件 (V) 遗漏了“单纯”。然而命题~\ref{prop:6-7-9} 只对 $v\in\mathcal M^*(A;J_z)$ 证明参数正则性，而定理~\ref{thm:6-7-11} 证明的步骤 4 又直接以 该命题推出 (V)；对任意多重覆盖球面要求同一结论一般过强，故依这两处上下文补正。（径改）
\item[§6.7，182-199] 原书在这里误把微分方向的切向量元组代入非线性映射 $\pi^I$；但 $\pi^I$ 的定义域是模空间中的基点，紧接着的 $w_i=z_i$ 或 $z_{\alpha_i}$ 也只描述基点， 故改用相应的 $(\widetilde{\mathbf u},\mathbf z)$。（径改）
\item[§6.7，182-199] 原书在式~\eqref{eq:6-7-16} 的第二个等式后只写 $\alpha,\beta\in T_i$，遗漏了 $\alpha E\beta$；节点 $z_{\alpha\beta}$ 仅沿树边定义，且该式 正是式~\eqref{eq:6-7-15} 沿边的化简，故补上边关系。（径改）
\item[§6.7，182-199] 原书只说 $V_i\subset T_i$ 由“通过一条边与 $S_i$ 相连”的所有顶点组成， 未排除 $S_i$ 自身；按字面会纳入 $S_i$ 内部彼此相邻的幽灵顶点，并与紧接着的 “$u_\alpha$ 对每个 $\alpha\in V_i$ 都非常值”以及映射 $f$ 的唯一性冲突，故补明 $V_i\subset T_i\setminus S_i$。原书还把随后映射的靶写成 $S_i$，但在 $du_{\beta_i}\not\equiv0$ 的情形又令 $S_i=\varnothing$、$f(\beta_i)=0$；为同时涵盖这两种 情形，正文把靶补为 $S_i\cup\{0\}$。（径改）
\item[§6.7，182-199] 原书在这里仅要求 $\beta\in T_i\setminus\{\alpha\}$ 时消失；但森林中还可能 有不属于 $T_i$、却位于同一纤维 $w_i$ 的分支，而式~\eqref{eq:6-7-17} 要求这些分支不受 $Y_{w_i}$ 扰动。前面两次运用同一构造时也都对所有满足 $z_\beta=w_i$ 的其他分支施加了 消失条件，故据此补全。（径改）
\item[§6.7，182-199] 原书在尚未定义 $\pi_{\alpha\beta}$ 时，称其与到 $S^2\times S^2$ 的投影之复合映射横截；但复合的类型与紧接着的原像定义唯一确定 $\pi_{\alpha\beta}=(\pi\times\pi)\circ\operatorname{ev}_{\alpha\beta}$，原文亦把复合次序写反。（径改）
\item[§6.7，182-199] 原书说紧化中的“每个稳定映射”都有至少一条边；但紧化包含原主层，而主层 对应无边树。此处伪循环论证所考察的是逃离主模空间的序列所产生的新增边界极限，只有这些 稳定映射必有至少一条边，故据此限定。（径改）
\end{description}

\subsection*{第\,7\,章}
\begin{description}[font=\normalfont,leftmargin=6em,itemsep=0.2em]
\item[§7.1，201-209] 原书把 $\overline{\mathcal M}_{0,I}$ 的最高维上同调群印作 $H^{\#I-3}$，并把 $\beta_{k,I}$ 的余维印作 $2\#I$；但 $\dim_{\mathbb R}\overline{\mathcal M}_{0,I}=2\#I-6$，且 $\beta_{k,I}\in H_{2(k-\#I)}(\overline{\mathcal M}_{0,k})$，故两处次数都唯一应为 $2\#I-6$。（径改）
\item[§7.1，201-209] 原书此处把比较对象印作 $\mu(A,T)$；但紧邻上文给出的 $T$-层预期维数 本身就是 $\mu(B,k)-2e(T)$（当 $B=A$ 时即为 $\mu(A,T)$），而且不同边界层的树 $T$ 并不固定。伪循环论证所需、并由 $e(T)>0$ 得到的统一比较对象，是主层维数 $\mu(A,k)$。（径改）
\item[§7.1，201-209] 原书此处漏掉了除以重参数化群 $G$ 所减少的六个实维数，把 $\mathcal M^*_{0,0}(A;J)$ 的预期维数印作 $2n+2c_1(A)$；由紧邻的商空间定义以及 $\mu(A,0)=2n+2c_1(A)-6$，唯一应补为这里的数值。（径改）
\item[§7.1，201-209] 原书说存在“唯一的二次齐次多项式 $f$”，但 $f$ 的任意非零常数倍定义 同一零点集；唯一的是 $f$ 的射影类，等价地说，是二次曲线 $C_f$。（径改）
\item[§7.2，210-220] 原书在条件 (b) 中把主层误写成 $\mathcal M^*_{0,k}(A;J)$，并在下文证明中重复了这一写法；但此处讨论的是 $\operatorname{ev}_{J_0}$ 的边界，且 $J$ 尚未指定，故两处均应为 $\mathcal M^*_{0,k}(A;J_0)$。（径改）
\item[§7.2，210-220] 原书把这个线性映射两端的一形式空间底空间都印作 $M$；但 $u:S^2\to M$，且紧邻定义的 $\Phi(u,\xi)$ 与 $\pi_{u,J}$ 都作用在 $S^2$ 上的形式， 故两端唯一应以 $S^2$ 为底。（径改）
\item[§7.2，210-220] 原书这里只要求 $z_1\ne1$、$z_2\ne0$，但这不足以保证前文所要求的 $w\notin\{0,1\}$ 以及 $z\notin\{0,1,\infty,w\}$。由显示公式直接化简，这些条件还要求 $z_1\ne0$、$z_2\ne1$、$z_1\ne z_2$；加上它们仍给出一个开稠密的 $x_4$ 集合，足以完成 一般点处的计数。（径改）
\item[§7.3，220-225] 原书称 $\widetilde J$ 与 $\widetilde\omega$ 相容；但本段只假设 每个 $J_z$ 被 $\omega$ 驯服，这只能推出乘积结构被 $\widetilde\omega$ 驯服。相容性还需 额外假设每个 $J_z$ 与 $\omega$ 相容。（径改）
\item[§7.4，225-239] 原书此处把求值映射的定义域印作 $\mathcal M^*_{0,1}(A;J_0)$；但上一句讨论的是包含给定非常值球面的完整模空间 $\mathcal M_{0,1}(A;J_0)$，命题结论也包括多重覆盖球面。故据此补去星号。（径改）
\item[§7.4，225-239] 原书把 $G=\operatorname{PSL}(n+1,\mathbb C)$ 称为“general linear group”；但该记号表示射影特殊 线性群，故据群的明确定义补正名称。（径改）
\item[§7.4，225-239] 原书把矩阵 $F(u)$ 的最大秩印作 $\max(d+1,n+1)$；但 $F(u)$ 是 $(n+1)\times(d+1)$ 矩阵，其秩至多为 $\min(d+1,n+1)$，故唯一补正。（径改）
\item[§7.4，225-239] 原书把范围印作 $d\ge0$；但上文 $N_d:=\mathrm{GW}^{\mathbb{CP}^2}_{dL,3d-1}$ 在 $d=0$ 时会有负的标记点数，因而 $N_0$ 未由该式定义，且下述证明从正次数开始。故唯一补为 $d\ge1$。（径改）
\item[§7.4，225-239] 原书在此直接以非齐次式 $z_0^2-z_1^3$ 标记 $\mathbb{CP}^2$ 中的曲线；这可理解为仿射图 $z_2=1$ 上的局部尖点 方程，若写成齐次方程则为 $z_0^2z_2-z_1^3=0$。（径改）
\item[§7.4，225-239] 原书此句末误印作“for every $k$”；本段只以 $j$ 标号节点，故据紧邻定义 唯一补为“对每个 $j$”。（径改）
\item[§7.4，225-239] 原书此处误称 $\mathcal W$ 的复维数为 $\delta$；它是 $\delta$ 个横截超曲面的交，故复余维为 $\delta$，这也与本步 陈述及下一步的维数计算一致。（径改）
\item[§7.4，225-239] 原书把多项式 $f=f_0\cdots f_k$ 印作属于射影化空间 $\mathcal P_d$； 但 $f_i$ 及其乘积都是 $\mathcal F$ 中的多项式，而下文才使用其射影类 $[f]\in\mathcal P_d$， 故唯一补为 $f\in\mathcal F_d$。（径改）
\item[§7.4，225-239] 原书此处的格罗莫夫--威滕伪循环定义域漏印星号；定理~\ref{thm:7-4-1} 及本节前文的同一求值映射均以单纯曲线模空间 $\mathcal M^*_{0,3d-1}$ 为定义域，故唯一补出。（径改）
\item[§7.4，225-239] 原书在命题~\ref{prop:7-4-9} 的两个证明中把该映射的定义域简写成 $\mathcal M(\mathbb{CP}^n,2L;J_0)$；但靶有五个因子，且后文的点为 $[u,\mathbf z]$，故必须使用带五个标记点的单纯曲线模空间 $\mathcal M^*_{0,5}(\mathbb{CP}^n,2L;J_0)$。（径改）
\item[§7.4，225-239] 这段概述按字面只适用于 $d=2$：当 $d>2$ 时，复余维为 $2$ 的 约束与 $\mathbb{CP}^d$ 的一般交仍是复余维为 $2$ 的子空间，而不是点。一般情形应以下文 对 $n$ 的逐步归纳为准。（径改）
\item[§7.4，225-239] 原书在本证明的三处求值映射中把定义域下标印作 $k$；但本段固定了 $3d-1$ 个约束，靶为 $(\mathbb{CP}^n)^{3d-1}$，且式~\eqref{eq:7-4-8} 也使用 $3d-1$ 个标记点，故三处统一补正为 $3d-1$。（径改）
\item[§7.4，225-239] 原书先定义投影中心为 $p_0$，但在该式及随后两个补集中误写成未定义的 $p$；由投影中心及本段末句再次出现的 $p_0$，唯一补为 $p_0$。（径改）
\item[§7.5，239-255] 原书把这里两个循环的维数印作 $2(k-2\#I)$；但由 $Y_{k,I}=\pi_{k,I}^{-1}(\mathbf w_I)$ 的余维以及本节前面对 $\beta_{k,I}$ 次数的定义，维数均为 $2(k-\#I)$，故据此补正。（径改）
\end{description}

\subsection*{第\,8\,章}
\begin{description}[font=\normalfont,leftmargin=6em,itemsep=0.2em]
\item[§8.1，257-264] 原书此处印作 ``nondegenerate''；但 $\pi^{-1}(\gamma)$ 是奇维超曲面， 而且紧接着的句子明确给出这一限制的一维核，故应理解为具有最大秩。（径改）
\item[§8.2，264-270] 原书把 $u_t$ 的定义域印作 $\Sigma$；但此处 $u_t^*\widetilde\omega$、$du_t$ 以及沿 $\gamma$ 的纤维间平行移动都要求 定义域为起始纤维 $M_{\gamma(0)}$，故据此补正。（径改）
\item[§8.2，264-270] 原书把严格不等式写成对每个竖直向量都成立，因而错误地包含了零向量； 驯服条件只对非零向量要求严格正，故补上这一限制。（径改）
\item[§8.2，264-270] 原书把第一行 $E^{\mathrm{Vert}}(\widetilde u;U_\alpha)$ 的积分域误印为 $\Sigma$；该量是 $U_\alpha$ 上的局部能量，且随后各行均在 $U_\alpha$ 上积分，故补正为 $U_\alpha$。（径改）
\item[§8.3，270-277] 原书式~(8.3.1) 将扰动项印作 $X_H(u)$，但本节没有定义未加波浪号的 $u$，且紧接着把该一形式定义为 $X_H(\widetilde u)$；据此补正自变量。（径改）
\item[§8.3，270-277] 原书把椭圆正则性提升后的空间印作 $W^{\ell,p}$；但 $\eta$ 起初属于 $L^q$，其中 $q$ 是 $p$ 的对偶指数，对伴随算子 $D^*_{H,\widetilde u}$ 应用椭圆正则性时积分指数仍为 $q$。这与第~\ref{sec:6-7} 节 同型论证中的指数一致。（径改）
\item[§8.4，277-279] 原书把纤维写成 $(M,\omega)$；但本例只定义了 $\omega_\lambda$，且构造 所用环路明确位于 $\Ham(M,\omega_\lambda)$，故纤维的辛形式应为 $\omega_\lambda$。（径改）
\item[§8.5，279-285] 原书此处及本节后文 多次把整族 $\{\widetilde A_\alpha\}$ 统称为“球面同调类”。但式~(8.5.2) 规定 $\pi_*\widetilde A_0=[\Sigma]$，且根分支 $\widetilde u_0:\Sigma\to\widetilde M$ 是截面；只有 $\alpha\ne0$ 的类满足 $\pi_*\widetilde A_\alpha=0$ 并由纤维中的球面表示。故后文均按根类为截面类、其余为 竖直球面类表述。（径改）
\item[§8.5，279-285] 原书在定义 8.5.2 与定理 8.5.3 中均把此处的 $\Sigma^I$ 印作 $\Sigma_I$；但前文只定义了由 $I$ 标号点组成的乘积 $\Sigma^I$，而 式~\eqref{eq:8-5-5} 及下文普适投影的靶也都是 $\Sigma^I$，故两处统一补正为上标。（径改）
\item[§8.5，279-285] 原书把此非负性写成对每个 $\alpha\in T$ 都成立； 但条件~(8.5.1) 只约束纤维中的球面类，而根类满足 $\pi_*\widetilde B_0=[\Sigma]$，是截面类，其竖直 Chern 数不必非负。故唯一可推出的范围是 $\alpha\ne0$。下文投影次数之和为一还蕴含根类重数为一，所以这一补正不影响随后所需的 维数比较。（径改）
\item[§8.6，285-294] 原书把定义域印作单个 $H_*(M)$，但紧接着的定义以 $k$ 个同调类 $\alpha_1,\ldots,\alpha_k$ 为输入，严格说这是一个 $k$ 重多线性映射，等价地定义在 $H_*(M)^{\otimes k}$ 上；式中再由交叉积把纯张量送入 $H_*(M^k)$。（径改）
\item[§8.6，285-294] 原书此处把同调群印作 $H_2(\widetilde M;\mathbb Z)$；但本例的总空间 从头至尾都是 $\widetilde M_f$，且紧接着也在其中考察全纯曲线，故补回下标 $f$。（径改）
\item[§8.6，285-294] 原书此处印作 $\widetilde\sigma_K\in H^1(\widetilde M_\psi)$；但紧接着要取 $d\widetilde\sigma_K$ 并用它改变联络形式，故该对象唯一应属于 $\Omega^1(\widetilde M_\psi)$。（径改）
\end{description}

\subsection*{第\,9\,章}
\begin{description}[font=\normalfont,leftmargin=6em,itemsep=0.2em]
\item[§9.1，295-309] 原书把该哈密顿函数的类型误印为 $\mathbb R\to M$；但输入是 $(t,x)$，且 $H_t(x)$ 为实值，故其源、靶唯一应为 $\mathbb R\times M\to\mathbb R$。（径改）
\item[§9.1，295-309] 原书此处把圆柱的第二因子误印为 $\mathbb R$；但 $u_T$ 的定义域及随后对 $t\in[0,1]$ 的积分都唯一确定该因子应为 $\mathbb R/\mathbb Z$。（径改）
\item[§9.1，295-309] 原书把 $\phi$ 称为到整个 $S^2$ 的微分同胚；但该指数映射的像是 $S^2\setminus\{0,\infty\}$，且前文也正是用它等同圆柱与双穿孔球面，故补正其靶。（径改）
\item[§9.1，295-309] 原书在本式及下文一般偶对的对应公式中两次误印为 $\operatorname{ev}_{J,H}$；本段扰动变量始终是 $K$，且紧邻处定义的也是 $\operatorname{ev}_{J,K}$，故据此补正。（径改）
\item[§9.1，295-309] 原书此处把圆盘记作 $B$，但本节只定义了单位圆盘 $\mathbb D$，且紧接着仍以 $e^{2\pi it}$ 参数化其边界，故统一改为 $\mathbb D$。（径改）
\item[§9.1，295-309] 原书在随后直接使用了未定义的 $J_T$。这里根据中央圆柱上必须还原 弗洛尔方程所用的 $J_t$，以及驯服殆复结构空间的 可缩性，补出这一光滑延拓族。（径改）
\item[§9.2，310-322] 原书此处把点 $x$ 误写成 $x\in T_xM$；由随后的切空间 $T_xM$ 以及“对每个边界点”的上下文，唯一应为 $x\in\partial M$。（径改）
\item[§9.2，310-322] 原书此处印作 $f(\phi_t(x))=e^t$，但同时规定 $f:M\to(-\infty,0]$ 且 $\partial M=f^{-1}(0)$；在 $t=0$ 时原式给出 $f=1$。唯一相容的 归一化是 $e^t-1$。（径改）
\item[§9.2，310-322] 原书在第二行微分号内把前一行的差 $\partial_tg\,ds-\partial_sg\,dt$ 误印成了和；保留负号时只有这里写成差，才能得到 $(\Delta g)\,ds\wedge dt$。（径改）
\item[§9.2，310-322] 原书把 $\omega_{\mathrm{can}}$ 的类型误印为 $\Omega^1(T^*L)$；它是 $2$-形式，且紧随其后的局部公式也是二次形式。（径改）
\item[§9.2，310-322] 原书此处写作 $r<r_0$，但上下文只定义了水平集参数 $c$，并未定义 $r_0$；相应刘维尔场指向的正是 $Q_c$ 内侧 $r<c$。（径改）
\item[§9.2，310-322] 原书写作“$X$ 上”的殆复结构，但本习题中并未定义流形 $X$；$J$ 的定义域 只能是带 $G$-作用的纤维 $M$。（径改）
\item[§9.2，310-322] 原书式~\eqref{eq:9-2-8} 把 $\mu(x)\in\mathfrak g^*$ 直接作为只对 $\xi\in\mathfrak g$ 定义的向量场 $X_\xi$ 的下标；除非另行指定一个 $\operatorname{Ad}$-不变内积来识别 $\mathfrak g^*\cong\mathfrak g$，该式没有类型意义。 原书未给出这种识别，故这里保留其最低解释度的记号。（径改）
\item[§9.2，310-322] 原书把靶空间印作本推论中未定义的 $M$；缩放公式及前后的哈密顿同痕均在 余切丛 $T^*L$ 中，故唯一应为 $T^*L$。（径改）
\item[§9.2，310-322] 原书把单位圆盘的条件误印为 $t^2+t^2\le1$；圆盘坐标为 $s+it$，故首项 唯一应为 $s^2$。（径改）
\item[§9.2，310-322] 原书把该函数的类型误印为 $\mathbb R\times[0,1]\to M$，却仍写三元输入 $(s,t,x)$ 与实值 $H_{s,t}(x)$；其唯一 相容的源靶是 $\mathbb R\times[0,1]\times M\to\mathbb R$。（径改）
\item[§9.2，310-322] 原书在同一页再次漏掉哈密顿函数的流形变量，把定义域印作 $\mathbb R\times[0,1]$；但 $H_t$ 是 $M$ 上的函数，且下文使用其哈密顿向量场，故补入 $M$。（径改）
\item[§9.3，323-329] 原书定义嵌入时只给出半径 $\delta$，随后直接使用 $U_r$，却漏写了保证 $\iota(z)$ 有定义所必需的 $0<r<\delta$。（径改）
\item[§9.3，323-329] 原书把势函数印作 $\log|z^2|$；紧接着的公式及齐次坐标中的标准 Fubini--Study 势函数均唯一确定此处应为 $\log|z|^2$。（径改）
\item[§9.3，323-329] 原书把识别映射印作 $h_\rho\circ\iota^{-1}$；按前一行给出的 $h_\rho$ 的源靶，这个复合映向欧氏球壳而非 $U_\varepsilon$，方向不合。此外，原书没有 定义这里的 $\varepsilon$。正文只保留证明所需的逆向识别，不擅自指定 $\varepsilon$ 的公式。（存疑照排）
\item[§9.4，329-344] 原书把二形式 $\psi^*\omega$ 与 $\omega_0$ 称为 ``homologous''；Moser 同痕所需且由前述积分确定的是二者的上同调类相同。（径改）
\item[§9.4，329-344] 原书在此印作 $U\times(S^2\setminus U)$；但当前曲线的第二坐标 $w$ 属于 $U$，其像落在 $S^2\times U$，而定理假设正是在这整个区域上给出标准形式。 原印的因子次序与区域不合。（径改）
\item[§9.4，329-344] 原书先只在 $1\le t\le T$ 定义 $f_t$，随后却把这句话的量词印作 ``all $t\ge1$''；此处唯一需要且有定义的范围是 $1\le t\le T$。（径改）
\item[§9.4，329-344] 原书把这里的环境空间印作 $\mathbb R^4$，但引理及整个 证明都在 $\mathbb R^{2n}$ 中进行。（径改）
\item[§9.5，344-353] 原书在本引理中把环境写成 $\mathbb R^{2n}$，但此处 $\Delta_r=D_r\times D_r$、与 $S^2\times S^2$ 的紧化以及紧随其后的整个证明都只针对 $\mathbb R^4$；故把环境唯一补正为 $\mathbb R^4$。（径改）
\item[§9.5，344-353] 原书在这一个包含式的左端把标准形式下标漏掉，写成了 $\operatorname{Symp}_V(\mathbb R^4,\omega)$；本证明始终固定的是 $\omega_0$，而缩放公式也保持 $\omega_0$，故补上该下标。（径改）
\item[§9.5，344-353] 原书把终点印作 $J_1=\psi_*J$，但此处只给定了初始标准结构 $J_0$，并且 随后要求 $u_1=\psi\circ u_0$；唯一相容的终点是 $\psi_*J_0$。（径改）
\item[§9.5，344-353] 原书同时固定所有 $u_t$ 的两个标记点之像，并要求 $u_1=\psi\circ u_0$；但一般的 $\psi$ 不固定 $[1:0:0]$ 与 $[0:1:0]$，故这些条件并不 相容。一种可能的补法是让标记点之像随 $t$ 移动，但原书未指定具体选择。后续实际需要的 是连接 $u_0$ 与 $\psi\circ u_0$ 的带标架曲线族，因而论证意图仍可辨认。（径改）
\item[§9.5，344-353] 原书此处写作 $\omega(\,\cdot,J_t\,\cdot)$，但本段的靶流形是 $\mathbb{CP}^2$，其固定辛形式为 $\omega_{\mathrm{FS}}$；本节此前的裸记号 $\omega$ 专指 $S^2\times S^2$ 上的乘积形式，故补全下标。（径改）
\item[§9.5，344-353] 原书把这里的参数范围写成 $\lambda\in\mathbb C$；但 $u_{\lambda\xi}$ 只对非零初速度 $\lambda\xi$ 定义， 故必须排除 $\lambda=0$。（径改）
\item[§9.6，353-359] 原书在这里把联络形式所属空间印作 $\mathcal T(\widetilde M_\psi)$，但本节只定义并使用 $\mathcal T_0(\widetilde M_\psi)$；紧随其后的公式也正是注~\ref{rem:8-2-11} 中该空间的 联络形式，故补上下标 $0$。（径改）
\item[§9.6，353-359] 原书把括号中“其他水平截面”说成对应 $H^k$ 的极小值；但前文已说明 $Z^k$ 正是极小值流形，而其余固定分支是极大值流形。故把 minima 唯一补正为 maxima。（径改）
\item[§9.6，353-359] 原书在此处求值映射的定义域把模空间类漏写了上标 $k$，印作 $\mathcal M(\widetilde A;J,\widetilde\omega)$；本段始终固定命题~\ref{prop:9-6-4} 的类 $\widetilde A^k$，故补全。（径改）
\item[§9.7，359-368] 原书右端印作 $\frac{c_t}{2}(d(r^2-1)\alpha)$：其一，括号使 $d$ 只施于标量 $r^2-1$，漏掉 $(r^2-1)d\alpha$ 项；其二，分子还漏了加项 $+1$（等价地系数少了 $1/2$）。事实上，左端在对角线上的限制为 $(c_t+1)\sigma$，而由 $-d\alpha=2\pi^*\sigma$， $d((r^2-1)\alpha)$ 在零截面上的限制为 $2\pi^*\sigma$，故右端唯一应为 $\frac{c_t+1}{2}d((r^2-1)\alpha)$。（径改）
\end{description}

\subsection*{第\,10\,章}
\begin{description}[font=\normalfont,leftmargin=6em,itemsep=0.2em]
\item[§10.2，373-376] 原书此处误印为 $u(0)=v(\infty)=1$；本例只定义了 $u^0,u^\infty$，故按紧邻定义唯一补正其上标与字母。（径改）
\item[§10.4，380-382] 原书在式~\eqref{eq:10-4-2} 最后一项中误印为 $2\pi c/\log(\varepsilon/\delta)$。由引理~\ref{lem:10-4-1}， $\|\nabla\beta\|_{L^2}=\sqrt{2\pi/\log(\varepsilon/\delta)}$，故该项唯一应为这里的 平方根形式。（径改）
\item[§10.5，382-391] 原书在这里把连通对空间缩写为 $\mathcal M(A^{0,\infty};J,c)$；但固定数据始终是偶对 $J^{0,\infty}=(J^0,J^\infty)$，且本节下一节也明确采用 $\mathcal M(A^{0,\infty};J^{0,\infty},c)$，故补出上标。（径改）
\item[§10.5，382-391] 原书在下式第一行把右端乘积的范数印作 $W^{1,p}$ 范数；但这里比较的是加权零阶范数与普通 $L^p$ 范数，且 引理~\ref{lem:10-4-2} 恰好只给出乘法算子 $W^{1,p}\to L^p$ 的估计。 若取 $W^{1,p}$ 范数则会无依据地要求截断函数的二阶导数，故补正为 $L^p$ 范数。（径改）
\item[§10.7，400-405] 原书在本段把关于满射性的结论写成第（iv）部分；但紧接着 的证明标题、四个步骤及证明结尾均明确指向定理~\ref{thm:10-1-2} 的第（v）部分，故补正。（径改）
\item[§10.8，406-412] 原书把这里的指标上界印作 $\ell^\infty$；但 $\mathbf w^0$ 是 $\ell^0$ 元组，且本节前后固定在 $0$ 侧的恰为前 $\ell^0$ 个标记点，故补正为 $\ell^0$。（径改）
\item[§10.8，406-412] 原书在右侧两个分支条件中把顶点下标印作 $\alpha_j$；但此处所固定的是 第 $\ell^0+j$ 个标记点，且同一分式的分子已经分别写成 $z_{\ell^0+j}$ 和 $z_{\alpha_{\ell^0+j}}$，故补正为 $\alpha_{\ell^0+j}$。（径改）
\item[§10.9，413-416] 原书取 $\rho=(\delta\delta_0)^2$，同时把 $B_\rho$ 定义为闭圆盘， 却在随后对其所有非零元素使用严格不等式 $|\lambda|<\rho$，从而不能覆盖边界。 定理只要求某个正半径，故改为任取 $0<\rho<(\delta\delta_0)^2$；于是闭圆盘上的 $R=|\lambda|^{-1/2}$ 也始终满足所需的严格下界。（径改）
\item[§10.9，413-416] 原书在这里定义元素时漏写输入 $\lambda$，印作 $\iota_{T,c}([\mathbf u,\mathbf z])$；本映射的定义域是 $\mathcal M_{0,T}^*(c)\times B_\rho$，而右端也显式依赖 $\lambda$，故补全。（径改）
\item[§10.9，413-416] 原书在这个光滑映射的定义域中把固定 $J$ 情形所需的 $\mathcal M^*(c)$ 漏掉星号，写成 $\mathcal M(c)$；同页单射性论证又把 $\mathcal M_{0,T}^*(c)$ 写成未定义的 $\mathcal M_{0,T}(c)$。由 注~\ref{rem:10-9-2} 及本证明开头粘合映射的定义域，两处均补上星号。（径改）
\item[§10.9，413-416] 原书在这个定义的右端再次写成 $\iota_{T,c,\beta}$，使公式成为自指； 这里是在原粘合映射 $\iota_{T,c}$ 的参数中作额外重标度，故补正右端下标。（径改）
\end{description}

\subsection*{第\,11\,章}
\begin{description}[font=\normalfont,leftmargin=6em,itemsep=0.2em]
\item[§11.1，417-435] 原书此处写作 “define an immersion”；但 very ample 的定义唯一给出到射影空间的嵌入， 故将 immersion 补正为 embedding。（径改）
\item[§11.1，417-435] 原书此处回指例 11.1.4(iii)； 但此处的双变量多项式环正是例 11.1.4(ii) 中群环在选定整数基后的表示， 而 (iii) 定义的是 Novikov 完备化，故改正例号。（径改）
\item[§11.1，417-435] 原书把 $\Lambda_\omega$ 回指为例 11.1.4(iv)；但 $\Lambda_\omega$ 明确定义于例 11.1.4(iii)，而 (iv) 定义的是普适 Novikov 环，故改正 例号。（径改）
\item[§11.1，417-435] 原书把 $a,b,c$ 中两因子拉回类之间的杯积都误印成加号；但紧接着的 格罗莫夫--威滕不变量乘积分解、多线性和次数条件都要求这里取外杯积，故改正。原书随后 的四点公式又直接使用了未定义的 $d_1,d_2,d$；按前三组变量的同一模式补出其定义。（径改）
\item[§11.2，436-441] 原书此处把符号误印为 $\varepsilon(\nu)=r(r+1)/2$，并在习题 11.2.3 中再次把 $\varepsilon(\alpha)$ 误印为 $r_\alpha(r_\alpha+1)/2$，两处都漏掉了 $(-1)$ 及幂号；前文已明确 $\varepsilon\in\{\pm1\}$，且紧接着的定义唯一确定了补正。（径改）
\item[§11.2，436-441] 原书把此式的因子印作 $(-1)^{\varepsilon(\nu,\alpha)}$，而紧接着的~\eqref{eq:11-2-8} 又把 $\varepsilon(\nu,\alpha)$ 定义为取值于 $\{\pm1\}$ 的符号，两式不相容。也可把 $\varepsilon$ 定义成模二奇偶数并保留外层 $(-1)$；本译稿依本节前后一贯的“符号”约定， 直接乘以 $\varepsilon(\nu,\alpha)$。（径改）
\item[§11.3，442-466] 原书在此定义 及紧接着的复化商中把环境空间误印为 $\mathbb C^n$；但本段的作用、坐标和矩映射均定义在 $\mathbb C^N$ 上，故补正为 $\mathbb C^N$。（径改）
\item[§11.3，442-466] 原书把左端下标的并集误印为 $I\cap I'$；由 $X_I=\{[x]\mid x_\nu=0\text{ 对 }\nu\in I\}$ 的定义，两个零坐标条件的交集对应 指标集的并集，紧接着的 $X_{\{\nu_1,\ldots,\nu_k\}}$ 也印证了这一点。（径改）
\item[§11.3，442-466] 原书最后一句把锥误写为 $\operatorname{cone}(I\setminus I_0)$；但 式~\eqref{eq:11-3-3}--\eqref{eq:11-3-7} 中始终出现的是 $I_0\setminus I$。原句的条件 $I\ne I_0$ 也须保留：这正说明除 $I=I_0$ 外不会产生乘法关系。（径改）
\item[§11.3，442-466] 原书把这个能级条件误印成逐坐标等式 $\pi|x_\nu|^2=-\kappa_\nu$。由 $\mu(x)=\pi\sum_\nu|x_\nu|^2w_\nu=\tau$ 与 $\sum_\nu\kappa_\nu w_\nu=-\tau$，正确且等价的条件是向量 $(\pi|x_\nu|^2+\kappa_\nu)_\nu$ 位于式~\eqref{eq:11-3-11} 的核中；原条件还会使紧接着 定义的 $\overline\mu$ 恒为零。（径改）
\item[§11.3，442-466] 原书在三角形的第三个不等式中误把 $\xi_1+\xi_2$ 印成 $\xi_1\xi_2$；由第三坐标 $-\xi_1-\xi_2\geq\kappa_3$ 唯一得到这里的和式。（径改）
\item[§11.3，442-466] 原书把这里的回指标号误印为 (11.3.7)。关系 (11.3.7) 是下式仍保留的 Stanley--Reisner 乘积关系；而 $\Delta_\nu\ne\varnothing$ 恰使单指标关系 (11.3.9) 不产生任何 $u_\nu=0$ 关系。（径改）
\item[§11.3，442-466] 原书把 $D$ 的环境格点误印为 $\mathbb Z^n$。由 $D=\dot\rho(\Lambda)$ 以及 $\dot\rho:\mathfrak t\to\mathbb R^N$ 可知， $D$ 的元素有 $N$ 个坐标，故应为 $D\subset\mathbb Z^N$。（径改）
\item[§11.3，442-466] 原书把 $\overline v(I)$ 误置于 $\overline{\mathfrak t}^{\,*}$；但每个 $\overline e_i$ 及扇锥 $\overline C_J$ 都位于 $\overline{\mathfrak t}$，故补正。（径改）
\item[§11.3，442-466] 原书把这些类的所在空间误印为 $H^2(M;\mathbb Z)$；本段对象始终是 $F_{n+1}$，故补正。（径改）
\item[§11.3，442-466] 原书称每个 $p_j$ 本身由 凯勒 形式代表，并称 $\sum_j\lambda_jp_j$ 在 $\lambda_j\geq0$ 且其和为正时恰为 凯勒，进而断言每个 $p_j(A)>0$。但下一段的对偶基满足 $p_j(A_i)=\delta_{ij}$，在 $n>1$ 时直接排除这些 断言。$p_j$ 是到相应格拉斯曼流形之投影的 凯勒 类拉回，故本身仅半正；凯勒 锥的 内部条件是所有 $\lambda_j>0$，而有效类上的结论相应为逐项非负、非零类至少一项为正。（径改）
\item[§11.3，442-466] 原书在这里把 $p_j$ 的定义域写成 $H_2(M;\mathbb C)$；本段所讨论的是 $T^*T_{\mathbb C}$，其余切纤维刚刚识别为 $H_2(F_{n+1};\mathbb C)$，故底空间应为 $F_{n+1}$。（径改）
\item[§11.3，442-466] 原书把点类生成元所在空间写成 $M$；这里从模空间、曲线到插入类始终 讨论的是 $G(k,n)$，故应为 $H^{2k(n-k)}(G(k,n))$。（径改）
\item[§11.3，442-466] 原书把这里的同调次数误印为 $2k$。零点集 $X\cong G(k,n-1)$ 的实维数为 $2k(n-k-1)=2k(n-k)-2k$，也正是次数 $2k$ 的类 $x_k$ 之庞加莱对偶的维数。（径改）
\item[§11.3，442-466] 原书在这里及下文再次定义 $c_t(E^*)$ 时都从 $i=1$ 起求和，却同时使用 $c_t(E^*)c_t(F^*)=1$、$1/c_t(E^*)$ 以及式~\eqref{eq:11-3-22} 右端常数项为一的乘积。 因此两个 Chern 多项式都唯一缺少 $c_0=1$；这里及式~\eqref{eq:11-3-22} 补出常数项。（径改）
\item[§11.3，442-466] 原书把闭包关系误印为 $b\preceq a$ 且 $b_i\leq a_i$。但由上面的精确跃迁 条件，闭包只能加入更特殊、余维更高的 Schubert 胞腔，即 $b_i\geq a_i$；例如 $G(1,2)$ 中 $a=1$ 的胞腔本身已是闭点。（径改）
\item[§11.3，442-466] 原书在这个等式中把联络 $A$ 的曲率 $F_A$ 误印为 $F_a$；前一句定义以及 紧接着的逆同构都唯一确定下标应为大写 $A$。（径改）
\item[§11.3，442-466] 原书在这句话中先把 $H_2(Z_k)$ 误印为 $H_2(M)$，随后又把 $\pi_2(Z_k)$ 漏写下标为 $\pi_2(Z)$；本段固定对象均为超曲面 $Z_k$，故统一补正。（径改）
\item[§11.3，442-466] 原书在此把不变量的曲线类印为 $dA$；但 $C$ 本身属于 $dA$，其 $m$ 重覆盖 属于 $mdA$，紧接着的图也明确属于 $A_0+mdA$，故补正。（径改）
\item[§11.3，442-466] 原书在这里把被扰动的乘积复结构误印为 $j\times J$；本段前文定义并始终 使用的是球面标准复结构 $i$，故补正为 $i\times J$。（径改）
\item[§11.3，442-466] 原书把右端印为 $\sum_{d\geq1}N_dq^d/(1-q^d)$，漏掉了覆盖次数的因子 $m^{-3}$。事实上，上一显示式给出 $\operatorname{GW}_{nA,3}^M=\sum_{d\mid n}N_dd^3$，而定义 $\operatorname{GW}_{nA,0}^M=n^{-3}\operatorname{GW}_{nA,3}^M$，故唯一得到这里的 $\sum_{d,m\geq1}N_dq^{md}/m^3$。（径改）
\item[§11.4，467-477] 原书把全上同调基的末指标印作 $n$；但同句中的 $n=\dim_{\mathbb R}M/2$，而本章固定的上同调基是 $e_0,\ldots,e_N$，故补正为 $N$。（径改）
\item[§11.4，467-477] 原书末项印作未定义的 $\mathcal S(a)*b$；本段定义的算子是 $\mathcal S(\phi)$，且左右模同态性质唯一要求此处为 $\mathcal S(\phi)(a)*b$。（径改）
\item[§11.4，467-477] 原书显示的映射不可能是所称同构：按紧随其后的定义，$\iota_!$ 把次数提高 二，而且投影公式给出 $\pi^*[S^2]\smile\iota_!b=0$。本证明实际由 $\iota_*$ 单射推出 $\iota^*$ 满射，从而在任选一个线性提升 $s:H^*(M)\to H^*(\widetilde M)$ 后得到 $a\otimes b\mapsto\pi^*a\smile s(b)$ 的非典范加性同构；原书未指定这样的提升，故此处 保留原显示式。（存疑照排）
\item[§11.4，467-477] 原书把这里的总空间印作 $\widetilde M_\psi$；但本步已取 $\psi_t\equiv\operatorname{id}$，式~\eqref{eq:11-4-6} 的左端、求和条件及类关系全都位于 $\widetilde M_\phi$，故补正下标为 $\phi$。（径改）
\item[§11.4，467-477] 原书把 $t$ 的指数印作本节从未定义的 $c_\nu$；紧前面的矩映射最大值公式及 Seidel 元首项均唯一给出该指数为 $\kappa_\nu$。（径改）
\item[§11.4，467-477] 原书把此处多项式环的末指标印作 $n$；但 本段由矩多面体的 $N$ 个面片定义变量 $u_1,\ldots,u_N$，前文环及 式~\eqref{eq:11-4-11} 也都以 $u_N$ 为末项，故补正为 $N$。（径改）
\item[§11.5，478-485] 原书在本注中把本节开头固定的基 $e_0,\ldots,e_N$ 的末指标三次改印成 $n$，从而写成 $\mathbb C^{n+1}$ 及求和至 $n$；但本章的 $n=\dim_{\mathbb R}M/2$，一般 并不等于 $\dim_{\mathbb C}\mathcal H-1$，故此处统一恢复为 $N$。（径改）
\item[§11.5，478-485] 原书第二行把逆度量展开式中的基向量误印为未定义的 $g_j$；这里唯一应为 $e_j$。（径改）
\item[§11.5，478-485] 原书把 $\mathbb C^m\times\mathcal H$ 的维数印作 $m+N$；但本节固定 $\mathcal H$ 的基为 $e_0,\ldots,e_N$，故其复维数为 $N+1$，总维数唯一应为 $m+N+1$。（径改）
\item[§11.5，478-485] 原书把 Lie 导数误印为 $\mathcal L_E(X):=-[X,Y]$，并把第二个恒等式首项 印成 $+[E,X*Y]$。按本书约定 $[X,Y]=dX\cdot Y-dY\cdot X$，前者唯一应为 $-[E,X]$；将这一导数作用于双线性乘法，后者首项也唯一应带负号。（径改）
\end{description}

\subsection*{第\,12\,章}
\begin{description}[font=\normalfont,leftmargin=6em,itemsep=0.2em]
\item[§12.1，487-499] 原书此处印作 ``semipositive monotone symplectic manifold''，同时要求“半正”和“单调”；但本节明确 讨论从单调情形到半正情形的推广，且后续命题与定理仅假设半正，故删去多余且与语境冲突的 ``monotone''。（径改）
\item[§12.1，487-499] 原书把流的初值印作 $\psi_0=0$；映射流的初值应为恒等映射。（径改）
\item[§12.1，487-499] 原书在能量恒等式中省略了波浪号及帽盘数据；但辛作用泛函只定义在带帽轨道的 覆盖空间上，而两端提升的相容性由延拓解 $u$ 拼接帽盘确定。（径改）
\item[§12.1，487-499] 原书此处把临界点集印作 $\operatorname{Crit}(H)$；本段已经固定的是 Morse 函数 $f$，且后续维数使用 $\operatorname{ind}_f(x)$，故改为 $\operatorname{Crit}(f)$。（径改）
\item[§12.1，487-499] 原书此处写作 $\Psi^\alpha$，但本段比较的是从弗洛尔上同调到量子上同调的 PSS 映射；按定理~\ref{thm:12-1-4} 及紧邻公式，该方向记为 $\Phi^\alpha$。（径改）
\item[§12.2，499-503] 原书从本段开头固定标记点 $w^\alpha,w^\beta,w^\gamma$，但随后在模空间条件 中改用未定义的 $z^\alpha,z^\beta,z^\gamma$，并把度量下标写成 $J_{z^i}$；这里统一恢复为 $w^\alpha,w^\beta,w^\gamma$ 与 $J_{w^i}$。（径改）
\item[§12.2，499-503] 原书把第二式的不稳定流形误印作 $W^u(x,H)$；此段只涉及 Morse 函数 $f$， 且右端正是 $f$ 的 Morse 指数，故改为 $W^u(x,f)$。（径改）
\item[§12.2，499-503] 原书在这里把三个 Morse--Witten 复形依次印作 $\operatorname{CF}^*(f^\beta;\Lambda_\omega)$、 $\operatorname{CF}^*(f^\gamma;\Lambda_\omega)$ 与 $C^*(f^\alpha;\Lambda_\omega)$；本节前文已把 Morse 复形统一记为 $\operatorname{CM}^*$，故将三处统一补正。（径改）
\item[§12.2，499-503] 原书此式的求和号仅列 $A^\beta+A^\gamma=A^\alpha$，因而把球面类 $A$ 与 临界点 $x^\beta,x^\gamma$ 留成自由变量，并在同调类约束中漏掉 $A$。紧邻的零维公式与 $\deg(x,A)=\operatorname{ind}_f(x)-2c_1(A)$ 唯一给出 $A+A^\alpha=A^\beta+A^\gamma$；这里也补齐全部求和变量。（径改）
\item[§12.3，503-505] 原书在这个公式右端把此前定义的链 $\eta$ 误印为 $\xi$；定义对象 以及紧随其后的配对相容式都唯一确定这里应为 $\eta(A\#\widetilde x)$。（径改）
\item[§12.4，505-514] 原书对任意 $\psi\in\operatorname{Symp}(M,\omega)$ 固定 $a$ 写出此式； 一般的自然性公式还应同时施加 $\psi$ 在量子上同调上诱导的作用。因此，原式固定 $a$ 的 版本至少还需假定 $a$ 在该作用下不变；特别地，对 $\psi\in\operatorname{Symp}_0(M,\omega)$（包括下文实际使用的哈密顿情形）这一条件成立。（径改）
\item[§12.4，505-514] 原书在中项把 $\iota^\kappa$ 的自变量误印为 $H$；但 $\Psi^{f,H_i}$ 的定义域是 $\operatorname{HF}_*(H_i)$，且本段使用的正则逼近列也是 $H_i$，故唯一应为 $\operatorname{im}\iota^\kappa(H_i)$。（径改）
\item[§12.5，514-518] 原书把这里的陪域印作尚未定义的 $\widetilde{\mathcal P}(H^\alpha)$；同页前后的双射与映射均明确给出陪域 $\widetilde{\mathcal P}(\psi^*H)$，故据此补正。（径改）
\item[§12.5，514-518] 原书在这里把 $\mathcal S^\alpha(\psi,\tau)$ 称为 Seidel 元；但前文已将它 定义为弗洛尔上同调上的 Seidel 同态，而本定理中的量子 Seidel 元是不带 $\alpha$ 的 $\mathcal S(\psi,\tau)$，故据此改正对象名称。（径改）
\item[§12.5，514-518] 原书把这里的量词印作 $A\in H_2(M)$；但 $A\#\widetilde x$ 只对第~\ref{sec:12-1} 节定义的甲板变换群 $H_2^S(M)$ 有定义，且本段 前后均使用该群，故补正。（径改）
\item[§12.7，524-529] 原书在这里把矩映射的陪域印作 $\mathfrak g$；但本节开头的定义以及同一 条件中的 $\tau\in Z(\mathfrak g^*)$ 均要求陪域为 $\mathfrak g^*$，故补正。（径改）
\end{description}

\subsection*{附录\,A}
\begin{description}[font=\normalfont,leftmargin=6em,itemsep=0.2em]
\item[§A.2，533-538] 原书第二版勘误（2019 年 11 月 27 日， \url{https://people.math.ethz.ch/~salamon/PREPRINTS/jholn-erratum.pdf}）指出， 印刷页 534--546 关于行列式线丛的讨论须重写以修正符号。以下保留锁定原书的论述； 涉及定向和同构的符号结论，不应视作已按勘误修正。（径改）
\item[§A.2，533-538] 原书称 $N$ 为正整数，随即又取 $N=\dim\operatorname{coker}D_0$；当 $D_0$ 满射时该维数为 $0$， 故这里据紧邻的构造改为“非负整数”。（径改）
\item[§A.2，533-538] 原书将下式左侧记作 $\mathcal F_1^0(X,Y)$，但上文已用同一记号定义 $k=1$ 的单一层；此处保留原记号，并以右侧并集明确其所指。 下文“简单交叉”处则按前一定义理解为核维数等于 $1$。（存疑照排）
\item[§A.2，533-538] 原书写作 $\ker D_t>0$， 将向量空间与数比较；据本页的分层定义改写为核维数大于零。（径改）
\item[§A.2，533-538] 原书只假设道路连续可微，却在下式及后续的核嵌入式写 $O(t^2)$； 单由该假设一般只能推出 $o(t)$。此处保留原式，未替原证明补加更强正则性假设。（存疑照排）
\item[§A.3，538-544] 原书下式右端的球心印作 $x_0$，但在平移映射 $\psi(x)=x+a$（其导数为恒等映射）下即不成立；证明中的归一化步骤及左端球心 均表明右端应为 $\psi(x_0)$，故据此补正。（径改）
\item[§A.3，538-544] 原书此处写“$\psi$ 在 $y_0$ 处可微”，但下一式及结论均讨论 $\psi^{-1}$ 的微分，且 $y_0$ 属于逆映射的定义域；据此补正主语。（径改）
\item[§A.3，538-544] 原书在此证明中直接使用常数 $c$，未先定义； 由紧接着的算子范数估计及末尾的球半径，可取任意 $c>0$ 且 $c\ge\|D^{-1}\|$。（径改）
\item[§A.3，538-544] 原书此处称所构造的 $\phi$ 为“smooth”，但本定理只假设 $f$ 为 $C^\ell$ 类，后面的证明也只得出 $\phi$ 为 $C^\ell$ 类； 故按本定理的正则性补正。（径改）
\item[§A.3，538-544] 原书此前明确取 $x_0\in\mathcal M=f^{-1}(0)$， 故在这段局部几何讨论中已有 $f(x_0)=0$。此句保留原文； 它只能理解为转入下一个命题的更一般设置，那里 $x_0$ 不必为零点。（存疑照排）
\item[§A.3，538-544] 原书此处右端印作 $\psi(B_{2\delta}(x_0;X))$，但 $\psi$ 未必在整个 $B_{2\delta}(x_0;X)$ 上有定义；由引理~\ref{lem:A-3-2} 的上界应为以 $\psi(x_0)$ 为中心、半径 $3\delta/2$ 的球。（径改）
\item[§A.3，538-544] 原书在此把 $dg(0)$ 写成 $DQ=\operatorname{id}_Y$，但按 $g(\eta)=f(x_1+Q\eta)$ 直接求导，应为 $dg(0)=df(x_1)Q$； $D=df(x_0)$ 仅在 $x_1=x_0$ 或另有条件时才与之相同。 后文迭代式可视为用冻结导数 $D$ 的迭代；这里保留原论证，不把它冒充精确的 Newton 步。（存疑照排）
\item[§A.3，538-544] 原书此处称转移映射“smooth”，但本定理只假设 $f$ 为 $C^\ell$ 类，且刚构造的坐标图及逆映射只保证 $C^\ell$； 故这里按结论所需的正则性表述。（径改）
\item[§A.4，545-547] 原书末段把 $f_0$ 的定义域印作 $U_0\cap\ker D$， 但定理及其证明定义 $f_0$ 于 $W=\psi(U_0)$；同段前一式也使用 $W\cap\ker D$， 故据此补正。原书此处又称其为 smooth 映射，在仅假设 $f$ 为 $C^\ell$ 时应为 $C^\ell$。（径改）
\item[§A.5，547-548] 原书此处只令 $x_0=0$，未说明还须将靶坐标平移 $-f(x_0)$；所引用的定理~\ref{thm:A-4-3} 以 $f(0)=0$ 为假设。 以下 $f$ 与 $y$ 均在平移后的坐标中理解。（径改）
\item[§A.5，547-548] 原书随后用到“$V$ 有界”，但此前只要求 $V\subset W$； 这里把局部邻域 $W$ 缩为有界集，使这一推理成立。（径改）
\item[§A.5，547-548] 原书从“对每个闭子集 $V_0\subset U_0$”直接过渡到最后的闭覆盖。 前面的闭性证明实际要求坐标逆像 $g^{-1}(V_0)$ 在整个 $X$ 中闭且有界； 仅知 $V_0$ 在 $U_0$ 中相对闭并不足够。一个可行的覆盖步骤是： 在各 $W$ 中选闭球 $B\subset W$，由可分性选出使 $g(B)$ 覆盖 $U$ 的可数族， 再直接对这些坐标闭球应用前面的闭性论证；原书未交代这一选择。（径改）
\end{description}

\subsection*{附录\,B}
\begin{description}[font=\normalfont,leftmargin=6em,itemsep=0.2em]
\item[§B.1，549-561] 原书差商括号内印为 $u(x-\xi/j)-u(x)$，其极限是沿 $\xi$ 方向导数的相反数，与紧随其后的弱导数恒等式符号相反；按该恒等式调整差商次序。（径改）
\item[§B.1，549-561] 原书在此未明言 $k\ge1$，但 $C^{k-1,1}$ 边界与证明中的 $k$ 项反射公式在 $k=0$ 时均无定义；命题应至少按 $k\ge1$ 理解。（径改）
\item[§B.1，549-561] 原书此处求和上界印为 $n$，但单位分解与下式均以 $j=0,\ldots,N$ 为指标；据此改为 $N$。（径改）
\item[§B.1，549-561] 原书权重印为 $k$，但对实值函数有 $a_{-k}=\overline{a_k}$，双边和恒为零，不能给出所述迹空间判据；据此改为 $|k|$。（径改）
\item[§B.1，549-561] 原书右端范数写作 $\|u\|_{W^{k,p}(\Omega)}$，但左端取 $u\circ\phi$，一般须在 $\phi(\Omega)$ 上控制 $u$；按复合映射定义域补正。原书 $\Phi(u\circ\phi)$ 依矩阵值 $\Phi$ 的类型在此理解为矩阵作用于向量，而非将 $u\circ\phi$ 代入 $\Phi$。（径改）
\item[§B.2，562-568] 原书平均值公式分母印作 $\omega_n r^2$；$\omega_n$ 是单位球面的 面积，$n$ 维球体积为 $\omega_n r^n/n$，故一般维数下应为 $\omega_n r^n$。原式仅在 $n=2$ 时成立。（径改）
\item[§B.2，562-568] 原书将 $K_{jk}=-nx_jx_k/(\omega_n|x|^{n+2})$ 写成 不加限制的二阶导数公式，随即另给 $K_{jj}$ 的正确公式；两式在 $j=k$ 时相矛盾，故前者须限于 $j\ne k$。（径改）
\item[§B.2，562-568] 原书只引用定理 B.2.7 和引理 B.2.2，但其基本解 式 (B.2.1) 未覆盖 $n=1$；这一维情形可直接由 $\Delta u=\partial_1^2u$ 得到。（径改）
\item[§B.2，562-568] 原书此处左端印作 $\partial_i u_j$，但本证明只定义了 标量函数 $u:=K_j*f$，未定义 $u_j$；依卷积求导式改为 $\partial_i u$。（径改）
\item[§B.2，562-568] 原书称这里所有内积都是 $L^2$ 内积；但当 $n=2$ 且 $\int f\ne0$ 时，$u=K_j*f$ 在无穷远处通常按 $|x|^{-1}$ 衰减， 不属于全局 $L^2$。此处 $\Delta u=\partial_jf$ 有紧支撑， 故 $\langle u,\Delta u\rangle$ 仍可作为积分使用。（径改）
\item[§B.2，562-568] 原书称此单个分量算子的范数为一；引理~\ref{lem:B-2-9} 仅给出上界一。实际上当 $j\ne k$ 时，其 Fourier 乘子为 $\xi_j\xi_k/|\xi|^2$，范数为 $1/2$；当 $j=k$ 时范数才为一。（径改）
\item[§B.2，562-568] 原书此步骤只要求 $h\in L^1$，但此时 $T$ 仅由 引理~\ref{lem:B-2-9} 在 $L^2$ 上定义，尚未给出一般 $L^1$ 函数的 $Th$ 的定义。步骤 3 的 $g$ 由 $f$ 在各立方体上取平均得到， 由 Jensen 不等式 $g\in L^2$，故实际使用的 $h=f-g$ 属于 $L^1\cap L^2$；将本步骤限于该交集足以完成证明。（径改）
\item[§B.3，568-571] 原书在一般 $n$ 下直接使用 $K$，但其定义式~\eqref{eq:B-2-1} 只覆盖 $n\ge2$。一维情形可由 $u''=f$ 直接处理；下文定理 B.3.2 的 卷积证明有同一省略。（径改）
\item[§B.3，568-571] 原书此处把 $L^p$ 收敛、$W^{k+2,p}$ 柯西性以及结论 均写在整个 $\Omega$ 上。但 $\Omega$ 未假设有界；例如在二维， $K*\beta f$ 可在无穷远处呈对数增长，不必属于 $L^p(\Omega)$。 证明所需的结论是在 $U$ 附近的局部正则性，故此处将这些断言限制在 相对紧的有界开集 $V$ 上。（径改）
\item[§B.3，568-571] 原书此处直接引用引理 B.2.2；该引理按其陈述要求 解与右端均有紧支撑，而这里的 $v=K*(\beta f)$ 一般没有紧支撑。 所需的正向分布恒等式可由该引理证明中的计算经逼近得到。（径改）
\item[§B.3，568-571] 原书在此估计及下文定理 B.3.2 中只写 $\overline{\Omega'}\subset\Omega$，未要求 $\Omega'$ 有界。 然而其证明要在 $\Omega'$ 外选支于 $\Omega$ 内的截断函数； 对一般无界 $\Omega'$，这种截断函数未必存在。这里补入 $\Omega'$ 有界这一充分条件。（径改）
\item[§B.3，568-571] 原书这条卷积估计的左端写为 $\|v\|_{W^{1,p}(\Omega)}$。当 $\Omega$ 无界时，全局范数仍可能不存在； 在证明中只须取有界的 $U$。（径改）
\item[§B.3，568-571] 原书此定理未要求 $\Omega'$ 有界； 此处补入该条件，以保证证明中可取紧支于 $\Omega$、 在 $\Omega'$ 的邻域上为 $1$ 的截断函数。（径改）
\item[§B.4，571-578] 原书未要求 $\Sigma$ 紧致，却给出全局 $W^{\ell,p}$ 结论； 下文证明实际得到的是 $W^{\ell,p}_{\mathrm{loc}}$。在非紧定义域上， 局部正则性不自动给出全局 索伯列夫范数有限；此处保留原断言。（存疑照排）
\item[§B.4，571-578] 原书仅要求 $\overline{\Omega'}\subset\Omega$，没有要求 $\Omega'$ 有界；若 $\Omega'$ 非有界， 这一包含关系并不保证统一的内缩距离，不能仅凭逐紧局部估计得到下文 对整个 $\Omega'$ 使用同一常数的估计。证明实际先在有界 Lipschitz 子域上估计； 命题~\ref{prop:B-4-9} 沿用相同开集条件，也有此疑点。此处保留原假设。（径改）
\item[§B.4，571-578] 原书将 $v_i$ 的陪域印作 $\mathbb R$，但下式要计算 $Rv_i$ 并与 $u:\Omega_0\to\mathbb R^{2n}$ 比较；据映射类型改为 $\mathbb R^{2n}$。（径改）
\item[§B.4，571-578] 原书此处极限迹写作 $s\mapsto u(s,t)$；由所用迹命题 及同句中的边界投影，改为 $s\mapsto u(s,0)$。（径改）
\item[§B.4，571-578] 原书沿用引理~\ref{lem:B-4-6} 的开集条件。局部正则性断言不受影响；但 (ii) 中对整个 $\Omega'$ 的 统一估计，不能仅从未限定有界性的闭包包含关系推出，参见前述来源脚注。（径改）
\item[§B.4，571-578] 原书在第一次提升后立即写 $J\circ u\in W^{2,p}_{\mathrm{loc}}$； 若仅有 $\ell=2$，推前的 $J$ 只有 $C^1$，该二阶复合正则性并不由此推出， 而且此时也无需再作下一步提升。这里按实际所需的归纳阶数表述。（径改）
\item[§B.4，571-578] 原书仅说“一致收敛”， 但假设只在逐个紧集上给出 $W^{1,p}$ 一致界；当 $\Sigma$ 非紧时， 此处由局部紧嵌入与对角子序列得到的是逐紧一致收敛。（径改）
\item[§B.4，571-578] 原书将它印作 $\mathbb R^{2n}\times\{0\}$，与 $u_\nu$ 的靶 $\mathbb R^{2n}$ 类型不合，且与紧邻的边界条件~\eqref{eq:B-4-8} 矛盾。（径改）
\end{description}

\subsection*{附录\,C}
\begin{description}[font=\normalfont,leftmargin=6em,itemsep=0.2em]
\item[§C.1，579-585] 原书将对偶丛的转换映射写为 $g_{\beta\alpha}^{T}$。若 $g_{\beta\alpha}$ 表示沿同一 $\alpha\to\beta$ 方向的纤维坐标变换，通常应为 $(g_{\beta\alpha}^{-1})^{T}$；原文未说明指标方向的约定，故此处保留原式而提示核对。（存疑照排）
\item[§C.1，579-585] 原书式 (C.1.8) 将右端的标量二形式误标为 $\Omega^{1,1}(\Sigma,E^*)$；由 $E^*$ 与 $E$ 的配对可知应属 $\Omega^{1,1}(\Sigma)$，此点亦见第二版 2019-11-27 勘误。（径改）
\item[§C.1，579-585] 原书对任意 $r>1$ 取其共轭指数 $q$，随即援引定理 C.1.10(i)；但该定理对 $k=1$ 还要求 $1-2/q\le\ell-2/p$，在 $\ell=1$ 且 $q>p$ 时不满足。此处保留原证明路线；未据此断定一般 $r$ 的结论错误，但该步尚需额外正则性论证。（存疑照排）
\item[§C.2，586-593] 原书上式在 $q=2$ 时的上界 $2q/(2-q)$ 分母为零。此时 $r=2s/(s-2)<\infty$；取充分接近 $2$ 的 $q'<2$，由 $W^{1,2}\hookrightarrow W^{1,q'}$ 及定理 B.1.12，仍得到 $W^{1,2}\hookrightarrow L^r$ 的紧性。（径改）
\item[§C.2，586-593] 原书此处的共同假设使用严格不等号，但定理 C.2.3 的陈述允许等号；本节后续步骤未单独讨论该端点。这里照录假设，端点情形的证明仍需核实。（存疑照排）
\item[§C.2，586-593] 原书步骤 5 此处第二个包含也误写成以 $\operatorname{coker}D_F^*$ 为目标；由伴随配对及本步骤下文可知，$\ker D_F^*$ 应映入 $\operatorname{coker}D_F$。（径改）
\item[§C.2，586-593] 原书步骤 6 此处将 $\xi$ 的目标丛误印为 $F$；$F$ 只定义在边界上，而式 (C.2.2) 与定理 (ii) 均要求 $\xi$ 是 $E$ 值截面。（径改）
\item[§C.2，586-593] 原书步骤 6 末句两处将作用于 $\eta$ 的算子误印为 $D_F$；其定义域是 $E$ 值截面，故此处据类型及式 (C.2.8) 改为 $D_F^*$。（径改）
\item[§C.2，586-593] 原书步骤 8 此处将作用于 $\eta-\eta_1$ 的算子误印为 $D_F$；该量是 $(0,1)$-形式，故据定义域及前一句改为 $D_F^*$。（径改）
\item[§C.2，586-593] 原书步骤 8 最后一次将作用于 $\eta,\eta_1$ 的算子误印为 $D_F$；据同一类型理由补为 $D_F^*$。（径改）
\item[§C.3，593-598] 原书把拉格朗日 Grassmann 流形还写成 $\operatorname{Sp}(2n)/\operatorname{GL}(n,\mathbb R)$。 这按通常的稳定子嵌入并非同一流形：当 $n=1$ 时该商为二维， 而 $L(1)=\operatorname U(1)/\operatorname O(1)$ 为一维。 辛群作用的完整稳定子还含对称矩阵的幂零部分；此处保留原式以待复核。（存疑照排）
\item[§C.4，598-604] 原书在定义 $A$ 之前便称可假定此等式成立； 一般的 $D$ 尚有下文定义的零阶项 $A$，该等式是以 $D-B$ 代替 $D$ 之后才得到的。此处依原书随后给出的构造调整叙述次序。（径改）
\item[§C.4，598-604] 原书实体页 617 在第一式的最后一项把 $\zeta_{12}(-s,t)$ 写成 $\zeta_{12}(s,t)$，与式 (C.4.6) 的反射自变量不符； 在第二式把系数导数 $(\partial_sC+i\partial_tC)$ 前误置负号。 两处按式 (C.4.6) 直接求导补正，未改变其余项。（径改）
\item[§C.4，598-604] 原书把反向曲面上的第二个边界马斯洛夫指标直接写成与第一个相同。 按 C.3 节从曲面诱导的边界定向定义，若第二份纤维仍取原复结构和同一平凡化， 其指标应变号；这里未明说为抵消该变号所需的共轭纤维约定。 此处保留原式，关于粘合后复线丛及陈数的约定需在复核时另行判定。（存疑照排）
\item[§C.4，598-604] 原书在闭曲面 $S$ 上仍写 $W_F^{1,p}(S,E)$，但此时并无边界条件 $F$；据定义域改为 $W^{1,p}(S,E)$。（径改）
\item[§C.5，604-613] 原书本步先称证明 $\psi\mapsto\psi^*i$ 的路径提升，却在随后的初值与方程中使用 $\psi_t^*j_t=i$，即 $\psi_t$ 为所需拉回映射的逆； 两者经微分同胚群中的取逆对应。此处保留原证明的方程。（存疑照排）
\item[§C.5，604-613] 原书将向量场空间写作 $\operatorname{Vect}(T\mathbb D)$， 但同一显示式将它等同于 $\Omega^0(\mathbb D,T\mathbb D)$； 按截面的定义，此处补为 $\operatorname{Vect}(\mathbb D)$。（径改）
\item[§C.5，604-613] 原书此处将外部点条件印为 $|\zeta|>0$，与紧前 $|\zeta|<1$ 时度为一的结论矛盾；下文也明确在 $|\zeta|>1$ 时使用零度结论。故改为 $|\zeta|>1$。（径改）
\item[§C.5，604-613] 原书将 $\hat u=du\cdot\xi$ 的零点条件转给 $\xi$ 时， 末项印作 $\xi(1)=1$；由线性切空间条件 $\hat u(1)=0$ 及 $du(1)$ 可逆， 应为 $\xi(1)=0$。下文核的计算亦使用此条件。（径改）
\item[§C.5，604-613] 原书在此由 B.4.2 直接说得到 $C^\ell$ 一致上界； 该定理陈述的是 $C^{\ell-1}$ 逐紧收敛子序列，足以推出下文所需的 $C^{\ell-1}$ 收敛。这里保留原文较强的中间断言。（存疑照排）
\item[§C.5，604-613] 原书以下上界式的分母写为 $|a_k|k^k$；若 $a_k=0$， 或在 $k=0$ 时不预设 $0^0$ 的约定，该表达式形式上未定义。 此处照录原式，不替作者增补约定。（存疑照排）
\item[§C.5，604-613] 原书下式令 $f_0(t)\equiv a_0$；若 $a_0=f(0)\ne0$，则按此式构造的 $f$ 在 $t\ge1$ 并不消失， 与习题开头的目标不符。此处照录原式，不替作者选定补救步骤。（存疑照排）
\item[§C.7，614-617] 原书在此将最后一个等号归于 (C.7.5)，但本引理只假设 $J$ 被 $\omega$ 驯服；(C.7.5) 要求 $J$ 与 $\omega$ 相容。此等号直接由无须相容假设的 (C.7.3) 得出。（径改）
\end{description}

\subsection*{附录\,D}
\begin{description}[font=\normalfont,leftmargin=6em,itemsep=0.2em]
\item[§D.1，619-622] 原书此处允许改用逆序列，但末尾仍用 $\phi^\nu=\psi^\nu\circ\rho^\nu$ 回推结论；若已改用逆序列，该等式便不再成立。需另行说明逆序列的逐紧一致收敛如何转回原序列，参见习题 D.1.3。（径改）
\item[§D.2，622-628] 原书此处写 $\alpha,\alpha'\in T$， 但这里要证明 $T_I$ 的连通性，且随后要求连接链落在 $T_I$；按紧前的 $T_I$ 定义，改为 $\alpha,\alpha'\in T_I$。（径改）
\item[§D.2，622-628] 原书接着声称链上每个顶点都可由四个已选指标之外的第五指标 $m$ 表出。若 $n=4$ 且树的两端分别带 $i,j$ 与 $k,\ell$，则 $[\alpha,\beta]$ 非空，却不存在这样的 $m$。此处保留原证明的断言；后续用这一表述证明树同态的步骤尚需补充论证。（存疑照排）
\item[§D.4，631-639] 原书在选边坐标时要求 $\alpha_n\in T_{\beta\alpha}$， 而此处又令 $I_{\beta\alpha}=\{1,\ldots,m\}$（$m<n$），从而 $n\notin I_{\beta\alpha}$。两处有向边的朝向约定不一致； 本段保留原记号，按此处两侧指标集合核对后续公式。（存疑照排）
\item[§D.4，631-639] 原书印刷第637页此式的左端及分式有误。此处将 $(1,m,m+1,n,i)$ 代入式~\eqref{eq:D-4-3}，并据下文 $w_{m,m+1,n,i}(w^0)=0$ 补正。第二版作者勘误所列分式为 $(w_{1,m,m+1,n}-1)/(w_{1,m,m+1,n}-w_{1,m,m+1,i})$； 它在 $w_{1,m,m+1,n}(w^0)=\infty$ 时趋于 $1$ 而非 $0$， 故与此处递推式和紧接的论证仍不符。（径改）
\item[§D.4，631-639] 原书此处写 $i,i'\in T_{\alpha\beta}$；$i,i'$ 是指标，而 $T_{\alpha\beta}$ 是顶点集。依本节 $I_{\alpha\beta}$ 的定义，改作指标所标记的顶点属于该分支。（径改）
\item[§D.5，640-642] 原书实体页 657 此分式的分母第二因子印为 $z_{\alpha i}-\zeta_{\alpha\ell}^\nu$。由前一行的交比 $w(z_{\alpha i},z_{\alpha j},z_{\alpha k},\zeta_{\alpha\ell}^\nu)$ 及 D.1 节的交比定义，该因子应为 $z_{\alpha k}-\zeta_{\alpha\ell}^\nu$，此处据此补正。（径改）
\item[§D.6，643-647] 原书实体页 658 此处写作 ``exactly three marked points, at least one of them nodal''；节点并非标记点。由本句后续的两种节点情形及稳定性定义，将 ``marked points'' 改作“特殊点”。（径改）
\item[§D.6，643-647] 原书实体页 659 写作 ``$M_{0,n+1}$ is the set of regular points''。这不能表示正则点集恰等于 $M_{0,n+1}$：例如 $n=3$ 时底空间 $\overline M_{0,3}$ 是单点，$\pi$ 的全部点均正则，而 $M_{0,4}$ 不含边界点。此处只将原句改为可成立的包含断言。（径改）
\item[§D.6，643-647] 原书实体页 660 在定义 $L_i$ 为截面 $Z_i$ 的余法丛拉回后，把纤维称为 ``tangent space''；余法丛的纤维应为纤维方向的余切空间。此处据紧邻定义补正。（径改）
\item[§D.6，643-647] 原书此步称限制类非零，但在所述范围的端点 $n=3$，$Z_{i,4}\cong\overline M_{0,3}$ 是单点，其 $H^2$ 为零；因此这一限制论证不能单独确定该端点的 $k$。此处保留原论证，提示端点须另行核验。（存疑照排）
\end{description}

\subsection*{附录\,E}
\begin{description}[font=\normalfont,leftmargin=6em,itemsep=0.2em]
\item[§E.1，653-658] 原书导言此句未写紧性或像不共享开集等限定；下文定理 E.1.2 对非单射点有限性要求定义域闭，定理 E.1.5 对两曲线还要求不存在像相同的非空开片。此处保留导言表述，精确适用条件以下文定理为准。（径改）
\item[§E.1，653-658] 原书此处定义 $\delta(u_0,u_1)$ 时将条件印作 $u(z_0)=u(z_1)$；由于 $z_i\in\Sigma_i$ 且 $u_i:\Sigma_i\to M$，据映射定义改作 $u_0(z_0)=u_1(z_1)$。（径改）
\item[§E.2，658-663] 原书此处写作“对每个 $k$ 次单位根 $\zeta$，$p(z)\not\equiv p(\zeta z)$”，漏排 $\zeta\ne1$；取 $\zeta=1$ 时该条件不可能成立。据前一定义及定理 E.1.2 的证明补上排除条件。（径改）
\item[§E.2，658-663] 原书在 $f_\varepsilon(z)=\prod_{j=0}^{k-1}(z+j\varepsilon)$ 之后将零点印作 $\{j\varepsilon\}$；由乘积中的正号可唯一确定为 $\{-j\varepsilon\}$，故据式 (E.2.1) 改正。（径改）
\item[§E.2，658-663] 原书此处把 $v$ 的靶印作 $\mathbb C^2$，但本引理设 $u:U\to\mathbb C^n$，式 (E.2.2) 也在 $\mathbb C^n$ 中定义 $v_\varepsilon$，故据引理陈述改作 $\mathbb C^n$。（径改）
\item[§E.2，658-663] 原书证明末句将上述两个定义域写作 $U_1\setminus u^{-1}(W)$ 与 $U_1\cap u^{-1}(W)$，而引理 (i) 要求在 $U$ 上成立；由于截断函数定义在 $U$ 且支集位于 $V'\subset u^{-1}(W)$，据构造改作 $U$。（径改）
\item[§E.2，658-663] 原书此处直接引用定理 E.1.2，但本命题允许 $\Sigma$ 非闭，而定理 E.1.2 的陈述要求 $\Sigma$ 闭；其证明中关于自交点孤立性的论证是局部的，不用闭性。此处据该证明改称“局部论证”。（径改）
\item[§E.2，658-663] 原书此处将 $\Psi_i(u(z_i))=0$ 与 $u(U_i)\subset W_i$ 分别印作 $\Psi_i(u_i(z_i))=0$、$u_i(U_i)\subset W_i$。$u_i$ 在下文才定义为坐标表达式，不能以 $z_i$ 或 $U_i$ 为自变量；据映射定义改正。（径改）
\item[§E.2，658-663] 原书此处以 $v\cdot v=u_\varepsilon\cdot u_\varepsilon$ 记互异参数点处局部交数之和。若把 $v\cdot v$ 解释为同调类的自交数，通常还须计入法丛的第一陈类项；参见下文定理 E.1.6 的证明。这里保留原式，将点积仅按双点交数和理解。（存疑照排）
\item[§E.2，658-663] 原书此处断言仅由互异参数点的局部交数之和等于同调类的自交数 $A\cdot A$。但本节定理 E.1.6 的证明在实体页 677 明确写出 $A\cdot A=\langle c_1(\nu_v),[\Sigma]\rangle+\sum\iota(v,v;z_0,z_1)$；一般还须计入法丛的 Euler／第一陈类项。此处保留原文断言，提醒读者核查。（存疑照排）
\item[§E.3，664-668] 原书实体页 682 只将 $\chi$ 定义在 $\mathbb D$ 上，随即对全部 $z\in\mathbb D$ 使用 $\chi(z/\varepsilon)$；当 $|z|>\varepsilon$ 时后者的自变量不一定在 $\mathbb D$。因 $\chi$ 在 $|z|\ge1/2$ 为零，按零延拓即可使该式有定义。（径改）
\item[§E.4，668-673] 原书实体页 684 仅把这条 Stokes 公式表述给实值 $f$，随后却代入复向量值函数 $\zeta^{-k}(\zeta-z)^{-1}v(\zeta)$；公式按实部、虚部逐分量线性延拓即可适用。（径改）
\item[§E.4，668-673] 原书实体页 685 把这里所需的积分绝对值上界印成 $\left|\int_{\partial\mathbb D_R}d\zeta/(\zeta-z)\right|$；该复积分恒为 $2\pi i$，不能直接上界带变权的边界积分。按紧接着的估计，实际所用的是绝对值置于被积式内的弧长积分，译文据此补正。（径改）
\item[§E.4，668-673] 原书实体页 686 仅以“$v$ 为 $C^2$ 函数”说明 $k=0$ 的极限为零；还须用假设 $v(z)=o(|z|^n)$（$n\ge1$）推出 $dv(0)=0$。此处补明所用条件。（径改）
\item[§E.5，673-678] 原书实体页 689 的逆映射导数公式以 $\det(\phi)$ 为分母，但 $\phi$ 是 $\mathbb R^2\to\mathbb R^2$ 的映射，不能取行列式；这里改为 $\det(d\phi)$。同一式中的第二个乘积下标 $\ell+j$ 改为与 $2|\alpha|$ 个指标一致的 $|\alpha|+j$。（径改）
\item[§E.5，673-678] 原书实体页 690 在此写 $\|h\|_{L^p(C_\varepsilon)}\le C\|h\|_{W^{1,p}(D_\varepsilon)}$，并把 $C$ 用作与 $\varepsilon$ 无关的常数。按通常未归一化的圆周测度，取常数函数 $h=1$ 即见此统一估计不成立；这里保留原书推导并提示该缩放疑点，结论可另由可去奇点定理核验。（径改）
\item[§E.5，673-678] 原书实体页 692 在 $E=z^{-k}R\in\mathcal O_k$ 后，把 $f=(\pi p)^{-1}\pi E$ 及 $(1+f)^{1/k}-1$ 均写作 $\mathcal O_{k+1}$；按本节乘积与复合规则，二者只能保证为 $\mathcal O_k$。乘上 $\psi_0\in\mathcal O_1$ 后，原书所需的 $\psi-\psi_0\in\mathcal O_{k+1}$ 仍成立。（径改）
\item[§E.5，673-678] 原书实体页 692 把 $F$ 写作 $u\circ\varphi-z^k=(\mathrm{id}-\pi)\circ u\circ\varphi$。这里 $u\circ\varphi$ 取值于 $\mathbb C^n$，而 $z^k$ 及 $\pi$ 取值于 $\mathbb C$；按已选 $a_0=(1,0,\ldots,0)$ 和 $\Pi=\ker\pi$，补全首坐标嵌入，使等式两侧均取值于 $\Pi$。（径改）
\item[§E.6，678-685] 原书实体页 694 写 $(0,\xi_0,x)\in\mathcal U$，但 $\mathcal F(u)=(\bar\partial_Ju,u(0),\partial_su(0))$，下一式又取 $\mathcal F^{-1}(0,x,\xi_0)$；故按此唯一参数顺序补正。（径改）
\item[§E.6，678-685] 原书实体页 695 由 $u_0=\mathcal F^{-1}(0,0,\xi_0)$ 直接写出 $u_0(z)=z\xi_0+z^2\bar\xi_0$。该多项式满足标准复结构的全纯方程，但在此仅假定 $J(0)=i$，一般非恒定 $J$ 下未必是 $J$-全纯解；这里保留原断言并标出疑点。（存疑照排）
\item[§E.6，678-685] 原书实体页 696 的引导段及实体页 698 的命题证明开头曾把 $H$ 写作 $F_1-F_2$，但紧接的引理 E.6.4、命题 E.6.1 的陈述及后续差分计算均取 $H=F_2-F_1$；据这些固定定义统一记号。（径改）
\item[§E.6，678-685] 原书实体页 696 把 $a,b,c$ 本身均写作属于 $\mathcal O_{k+1}$，却同时要求 $a(0)=c(0)=1$；二者不相容。由本引理证明中的 $a=1-\operatorname{Re}B_2$、$2b=\operatorname{Re}(A_2+iB_2)$、$c=1+\operatorname{Im}A_2$ 及 $A_2,B_2\in\mathcal O_{k+1}$，唯一补正是 $a-1,b,c-1\in\mathcal O_{k+1}$。（径改）
\item[§E.6，678-685] 原书实体页 697 此式积分上限印作 $\tau$，但 $\alpha$ 在紧邻定义为从 $0$ 到 $1$ 的积分；据定义改为上限 $1$。（径改）
\item[§E.7，686-694] 原书此处的最大值在 $S=\{f\}$ 时取自空集，未给约定；在该情形下，依定义直接有 $S_m^f=\{f\}$ 与 $A_m^f=f$。（径改）
\item[§E.7，686-694] 原书在此使用 $S_m$ 有限，但引理 E.7.2 的陈述只假设 $S_m$ 为 $k$-不变、$m$-分离集。有限性或足以取得共同定义域、统一分离常数的一致性条件尚需补入；后文具体应用中的 $S_m$ 确实有限。（径改）
\item[§E.7，686-694] 原书把上述估计的幂次写为 $m+1-|\alpha|$，但引理只假设 $S_m\subset\mathcal O_{k+1}$，不能推出 $m+1$ 阶的估计（取 $k=1$、$m>1$、$S_m=\{z^2\}$ 即是反例）。此处及下方 $\delta$ 的条件、式 (E.7.5) 的相应幂次据假设改为 $k+1$ 阶；结论不变。（径改）
\item[§E.7，686-694] 原书此处写 $z^2=x$；本引理固定任意 $k\ge1$，且条件 (E.7.4) 与下文各处均为 $z^k=x$，故补正。（径改）
\item[§E.7，686-694] 原书以下最小值表达式在 $x=0$ 时含分母 $|z|^m=0$，未说明约定；$x=0,y\ne0$ 时可由附近 $x\ne0$ 的零值邻域延拓理解，零点则另行处理。（径改）
\item[§E.7，686-694] 原书此处把 $y_0$ 所在空间印作 $\mathbb C^n$，但本引理的映射定义域及 $y$ 的维数均为 $\mathbb C^{n-1}$，据此补正。（径改）
\item[§E.7，686-694] 原书这里只给出 $|x|^{-1}$ 一项；这依赖于前述未被假设保证的 $m+1$ 阶估计。改用 $f\in\mathcal O_{k+1}$ 后还需计入第二项，乘上 $\Gamma(f)$ 的阶数后仍得到式 (E.7.8)。（径改）
\item[§E.7，686-694] 原书此处将后一个不等号印为 $m>m(f,g)$；由 $S_m^f=\{g\mid g-f\in\mathcal O_{m+1}\}$ 及 $m(f,f)=\infty$，应为 $m<m(f,g)$。这也正是相应差在 $m$ 阶消失的条件。（径改）
\item[§E.7，686-694] 原书此处写 $c_{\nu(j)}$，但本段只引入 $c_j$ 与阶数 $k_j$，未说明 $\nu(j)$ 的含义。这里保留原记号；首项系数须落在 $T_j$ 中，如何同 $c_j$ 精确对应，原文未说明。（存疑照排）
\item[§E.7，686-694] 原书此处对所有 $m\ge k$ 同时称 $a_m^f$ 已由引理 E.7.3 定义，但该引理 (v) 只在 $m>k$ 时给出此系数；下文 $m=k$ 的定义只用 $A_k^f$。（径改）
\item[§E.7，686-694] 原书只对 $f\ne g$ 的集合取最大值；当 $S$ 只有一个元素时该集合为空，$m_0$ 未定义。并入 $\{k\}$ 可涵盖此端点，且不改变 $|S|\ge2$ 时的取值。（径改）
\item[§E.7，686-694] 原书下一式第一个求和的下限印为 $m=k+1$，漏掉 $E_{\Gamma_k}=A_k^f$；依 $\Psi$ 的定义及紧邻上式，须从 $m=k$ 起求和。（径改）
\end{description}
